% =============================================================================
%  v13_review.tex -- review of thesis v13: corrections and ready-to-paste inserts
% -----------------------------------------------------------------------------
%  Page numbers refer to Master_s_thesis_13.pdf.
%  Text with a bar in the left margin is thesis text: copy it into the thesis.
%  Everything else is commentary.
%
%  The inserts refer to results by label. Add these \label{...} to the thesis
%  (or replace them by your own labels):
%     Assumption 3  ass:chang      Lemma 4        lem:bias
%     Assumption 5  ass:rate       Proposition 6  prop:sparse
%     Proposition 7 prop:ratefail  Proposition 8  prop:ratedr
%     Proposition 9 prop:stein     Lemma 10       lem:invariance
%     Lemma 11      lem:weights    Proposition 12 prop:decomp
%     Proposition 13 prop:pool     Corollary 14   cor:reference
%     Lemma 25      lem:decomp     Lemma 26       lem:variance
%     Lemma 27      lem:consist
%     Sections 4.1, 4.2, 5, 6, 7: sec:nonsparse, sec:basis, sec:mc, sec:emp, sec:conclusion
%     Tables: tab:d1, tab:pooling, tab:d23, tab:d23formula, tab:d4a, tab:d4aN, tab:d4b, tab:e1
%
%  Build:  cd review && pdflatex v13_review && biber v13_review && pdflatex v13_review (twice)
% =============================================================================
\documentclass[11pt]{article}
\usepackage[a4paper,margin=2.4cm]{geometry}
\usepackage[T1]{fontenc}
\usepackage[utf8]{inputenc}
\usepackage{lmodern}
\usepackage{amsmath,amssymb,amsthm,graphicx,enumitem,xcolor,booktabs,framed}
\usepackage[hidelinks]{hyperref}
\usepackage[style=authoryear,backend=biber,maxcitenames=2,uniquelist=false,uniquename=false]{biblatex}
\addbibresource{v13_refs.bib}
\providecommand{\E}{\mathbb{E}}
\newcommand{\ph}[1]{\textcolor{orange!85!black}{\textbf{[#1]}}}
\newcommand{\where}[1]{\par\medskip\noindent{\small\textcolor{black!60}{\textsf{\textbf{Where:} #1}}}\par\nopagebreak\smallskip}
\newcommand{\note}[1]{\par\noindent{\small\textcolor{blue!55!black}{\textsf{\textbf{Note:} #1}}}\par\smallskip}
\newcommand{\wrong}[1]{\textcolor{red!70!black}{#1}}
% thesis text: bar in the left margin (tables sit between two such blocks)
\newenvironment{ins}{\begin{leftbar}}{\end{leftbar}}
\theoremstyle{plain}
\newtheorem{lemma}{Lemma}
\renewcommand{\thelemma}{\unskip}
\newtheorem{proposition}[lemma]{Proposition}
\newtheorem{corollary}[lemma]{Corollary}
\newtheorem{definition}[lemma]{Definition}
\theoremstyle{remark}
\newtheorem{remark}[lemma]{Remark}
% --- thesis v13 numbers, for this preview only --------------------------------
\newlabel{ass:chang}{{3}{}{}{}{}}
\newlabel{lem:bias}{{4}{}{}{}{}}
\newlabel{ass:rate}{{5}{}{}{}{}}
\newlabel{prop:sparse}{{6}{}{}{}{}}
\newlabel{prop:ratefail}{{7}{}{}{}{}}
\newlabel{prop:ratedr}{{8}{}{}{}{}}
\newlabel{prop:stein}{{9}{}{}{}{}}
\newlabel{lem:invariance}{{10}{}{}{}{}}
\newlabel{lem:weights}{{11}{}{}{}{}}
\newlabel{prop:decomp}{{12}{}{}{}{}}
\newlabel{prop:pool}{{13}{}{}{}{}}
\newlabel{cor:reference}{{14}{}{}{}{}}
\newlabel{lem:decomp}{{25}{}{}{}{}}
\newlabel{lem:variance}{{26}{}{}{}{}}
\newlabel{lem:consist}{{27}{}{}{}{}}
\newlabel{sec:nonsparse}{{4.1}{}{}{}{}}
\newlabel{sec:basis}{{4.2}{}{}{}{}}
\newlabel{sec:mc}{{5}{}{}{}{}}
\newlabel{sec:emp}{{6}{}{}{}{}}
\newlabel{sec:conclusion}{{7}{}{}{}{}}
\newlabel{tab:d1}{{1}{}{}{}{}}
\newlabel{tab:e1}{{9}{}{}{}{}}
\newlabel{eq:sparsity1}{{7}{}{}{}{}}
\newlabel{eq:sparsity2}{{8}{}{}{}{}}
\setlist{itemsep=2pt,topsep=3pt}
\setcounter{tocdepth}{2}
\setcounter{table}{1}
\title{Review of thesis v13\\[2pt]\large corrections and ready-to-paste inserts}
\date{14 November 2026 (Section 6 revised after recreating Chang's estimate)}
\begin{document}
\sloppy
\maketitle
\vspace{-2.5em}

\section*{How to read this file}
Page numbers refer to \texttt{Master\_s\_thesis\_13.pdf}. I read all 44 pages, as
text and as rendered pages, redid every derivation in Sections 3--4 and in the
appendix by hand, recomputed every number in Section~5 from the CSV files in the
repository, and checked the citations I could against the papers themselves.
\begin{itemize}
\item Text with a bar in the left margin is thesis text. Copy it from
  \texttt{review/v13\_review.tex}; a grey line above it says where it goes.
\item \ph{Orange brackets} are numbers you still have to fill in from your own runs.
\item Blue notes are instructions to you, not thesis text.
\item ``Part~3'' means Section~3 of this file; ``Section~3'' always means the
  thesis.
\item The inserts use \verb|\ref{...}| with the labels listed at the top of the
  \texttt{.tex} file. In this PDF they print with the v13 numbers.
\item The citation keys are mine (file \texttt{review/v13\_refs.bib}). For
  entries you already have, replace my key by yours.
\item The inserts assume \verb|\providecommand{\E}{\mathbb{E}}|, \texttt{biblatex}
  (\verb|\textcite|, \verb|\parencite|) and, for the tables, \texttt{booktabs}.
\end{itemize}
Everything here supersedes the earlier files; you do not need
\texttt{thesis\_text.pdf} or \texttt{v12\_review.md} any more.

\tableofcontents

% =============================================================================
\section{Verdict and order of work}

\subsection{Where the thesis stands}
\begin{itemize}
\item \textbf{The theory holds.} I checked every step of Lemma~4,
  Propositions~6--9, Lemma~10--11, Propositions~12--13, Corollary~14 and
  Lemmas~25--27, including the proof of Proposition~7. The results are correct.
  What is wrong is in the write-up: about ten places where the printed statement
  or a proof step is not what you mean (Part~\ref{sec:math} below). The two
  that matter most are the missing conditioning on $X$ in the identifying
  assumptions, and $\hat p_k$, which is still defined in two different ways.
\item \textbf{The simulations are done and they confirm the theory.} D1 is
  written up. For D2/D3 and D4 the results are in the repository, but the thesis
  has only headings, or my chat notes pasted in. Part~\ref{sec:sims} gives
  the full text and the tables, with your numbers.
\item \textbf{Two parts are missing:} Section~6 (the application) and the
  conclusion. For Section~6 the baseline estimates now exist, and Chang's own
  estimate has been recreated (Part~\ref{sec:emp-ins}); the runs across reference
  groups remain.
\item \textbf{The introduction undersells the thesis.} It does not mention the
  pooling result (Proposition~13 and Corollary~14), which is the most original
  part, and it presents Propositions~6 and~8 as new, although they are known
  results applied to DML-DiD. See Part~\ref{sec:obvious}.
\end{itemize}

\subsection{Order of work}
\begin{enumerate}
\item Fix the mathematics (Part~\ref{sec:math}). Mostly pasting. Half a day.
\item Paste D2/D3 and D4 with their tables (Part~\ref{sec:sims}). Half a day.
\item Finish Section~6 (Part~\ref{sec:emp-ins}): the baseline and the recreation
  of Chang's estimate are done; the runs across reference groups remain.
\item Paste and adjust the conclusion (Part~\ref{sec:concl-ins}).
\item Framing, literature and bibliography (Parts~\ref{sec:front},
  \ref{sec:lit}, \ref{sec:obvious}).
\item Cut the appendix (Part~\ref{sec:obvious}) and do the notation and typo
  pass (Parts~\ref{sec:coherence}, \ref{sec:form}).
\end{enumerate}

% =============================================================================
\section{Corrections to the mathematics}\label{sec:math}

Items marked \textbf{(wrong as printed)} change the meaning of a statement or a
proof. The others are notation or wording.

\subsection{Section 3: Chang's framework (pp.~5--10)}

\paragraph{M1 (p.~6). The conditioning on $X_i$ is missing (wrong as printed).}
RO(ii) prints ``$\Pr(D_i=1_i)<1-\varepsilon$'', RO(iii) prints
``$\E[Y_i^0(1)-Y_i^0(0)_i,D_i=1]$'', and equation~(2) and the sentence after it
print ``$P(D_i=1_i)$''. In each case ``$\mid X$'' has dropped out and only the
subscript $i$ is left, so the identifying assumptions do not say what they
should. The same has happened in Assumption~3: ``$\E[U]=0$ and
$\E[V_1,D=0]=0$'', and in (c) and (d). Search the source for the macro or
replace-all that removed ``\verb+|X+''. The assumptions and (2) should read:
\begin{ins}
\begin{itemize}
\item[RO(ii)] \emph{Common support:} $\Pr(D_i=1)>0$ and
  $\Pr(D_i=1\mid X_i)\le1-\varepsilon$ almost surely, for some $\varepsilon>0$.
\item[RO(iii)] \emph{Conditional parallel trends:}
  $\E[Y_i^0(1)-Y_i^0(0)\mid X_i,D_i=1]=\E[Y_i^0(1)-Y_i^0(0)\mid X_i,D_i=0]$.
\end{itemize}
\[
  \theta_0=\E\left[\frac{Y_i(1)-Y_i(0)}{P(D_i=1)}\cdot
  \frac{D_i-P(D_i=1\mid X_i)}{1-P(D_i=1\mid X_i)}\right]
\]
\end{ins}
In Assumption~3: $\E[U\mid X]=0$ and $\E[V_1\mid X,D=0]=0$; (c)
$\E[U^2\mid X]\le C$; (d) $\E[V_1^2\mid X]\le C$. These are Chang's conditions,
and the proof of Lemma~27 uses (d) in the conditional form.

\paragraph{M2 (pp.~6--7). Small things in Section 3.1.}
\begin{itemize}
\item The assumption for repeated cross-sections is labelled ``RPC(i)'' and the
  estimand $\theta_0^{RPC}$. It should be RCS.
\item In (3) and (5), write $\hat g(X_i)$, and
  ``where $\hat p$ estimates $p_0\equiv P(D=1)$'' instead of ``$\hat p$ is the
  calculated constant $p_0$''.
\item The paragraph after (2) describes the untreated twice (``For the
  untreated, whose values \dots\ Conversely for the untreated''). Shorter and
  correct: ``A treated unit has weight $1/p_0$. An untreated unit has weight
  $-g_0(X_i)/\{p_0(1-g_0(X_i))\}$, the odds of treatment: untreated units that
  resemble the treated count more, and untreated units that do not resemble them
  count less.''
\end{itemize}

\paragraph{M3 (p.~7). The paragraph on Neyman orthogonality is broken.}
\begin{itemize}
\item ``repeated cross-sections (9)'' and, further down, ``means that
  \wrong{repeated outcomes} are locally insensitive'' are \verb|\ref|s that
  resolve to the wrong thing.
\item The sentence ``Chang proves for $\mathcal T_N\subset\mathcal T$ \dots\ with
  high probability that'' ends in Definition~1.
\item $D_r$ maps into $\mathbb R$, not $\mathbb R^s$ ($\theta_0$ is a scalar, and
  $s$ is the sparsity index later).
\item In the display for $D_0$, a bracket is missing and the direction
  $[\eta-\eta_0]$ is not shown. In Definition~1, $\rho_0$ should be $p_0$.
\item ``(Chang (2020); Chernozhukov et al. (2018).'' has unbalanced brackets.
\end{itemize}
Replacement for the text from ``The reason for introducing the adjustment
terms'' to the end of Definition~1:
\begin{ins}
The adjustment terms make the scores Neyman orthogonal
\parencite{chang2020double,chernozhukov2018double}: the moment conditions are
locally insensitive to the nuisance functions, so that errors in $\hat g$ and
$\hat\ell$ affect the estimator only through their product.

Formally, let $\eta=(g,\ell)\in\mathcal T$, where $\mathcal T$ is a convex subset
of a normed vector space, and for $\eta\in\mathcal T$ and $r\in[0,1)$ define the
directional (Gateaux) derivative
\[
  D_r[\eta-\eta_0]\equiv\partial_r\,\E\big[\psi\big(W,\theta_0,p_0,\eta_0+r(\eta-\eta_0)\big)\big],
  \qquad
  \partial_\eta\E\,\psi(W,\theta_0,p_0,\eta_0)[\eta-\eta_0]\equiv D_0[\eta-\eta_0].
\]
Let $\mathcal T_N\subset\mathcal T$ be a nuisance realisation set, i.e.\ a set
that contains the estimator of $\eta_0$ with probability approaching one.

\begin{definition}
The score $\psi$ obeys the Neyman orthogonality condition at
$(\theta_0,p_0,\eta_0)$ with respect to $\mathcal T_N$ if $D_r[\eta-\eta_0]$
exists for all $r\in[0,1)$ and $\eta\in\mathcal T_N$, and
$\partial_\eta\E\,\psi(W,\theta_0,p_0,\eta_0)[\eta-\eta_0]=0$ for all
$\eta\in\mathcal T_N$.
\end{definition}

\textcite[Lemma~3.1]{chang2020double} shows that $\psi_1$ and $\psi_2$ obey this
condition. Note that the scores are orthogonal with respect to the functions
$(g,\ell)$ only, not with respect to the constants $p_0$ and $\lambda_0$. The
estimation error in $\hat p$ and $\hat\lambda$ therefore does not vanish at first
order. It does not create bias, but it enters the asymptotic variance, which is
why the variance estimators below contain the terms $\hat G_p(D-\hat p_k)$ and
$\hat G_\lambda(T-\hat\lambda_k)$.
\end{ins}

\paragraph{M4 (p.~8, p.~28, p.~33). $\hat p_k$ is defined in two ways (wrong as printed).}
Definition~2 computes $\hat p_k$ on the auxiliary sample,
$(N-n)^{-1}\sum_{i\in I_k^c}D_i$. The appendix (p.~28 and p.~33) computes it on
the fold. Footnote~2 (p.~33) says ``on the auxiliary sample'', directly under a
sentence that says the fold. The proof of Lemma~25 says ``the mean of $N-n$
Bernoulli variables''. Your code uses the fold, and so does Chang's.
\begin{itemize}
\item Definition~2, sentence after the two displays:
\begin{ins}
where $\hat p_k=n^{-1}\sum_{i\in I_k}D_i$ and
$\hat\lambda_k=(N-n)^{-1}\sum_{i\in I_k^c}T_i$, and
$(\hat g_k,\hat\ell_{1k},\hat\ell_{2k})$ are estimators of the nuisance functions
constructed using the auxiliary sample.
\end{ins}
\item Footnote~2 (move it to Definition~2):
\begin{ins}
\textcite[Definition~3.1]{chang2020double} writes
$\hat p_k=n^{-1}\sum_{i\in I_k^c}D_i$, which sums over the auxiliary sample but
divides by the size of the fold. I compute $\hat p_k$ on the fold $I_k$, as his
replication code and his proofs do. With $\hat p_k$ from the auxiliary sample,
the numerator and denominator of $\tilde\theta_k$ would be independent, and by
Jensen's inequality $\tilde\theta_k$ would have a bias of about
$\theta_0(1-p_0)/\{p_0(N-n)\}$. In Chang's simulation design with $K=5$ this is
0.019, which is of the same order as the standard deviation of the estimator
(0.025). With $\hat p_k$ from the fold, the
treated part of $\tilde\theta_k$ is an average over the treated units of the
fold, and this bias does not arise. $\hat\lambda_k$ is computed on $I_k^c$,
because the outcome nuisance is fitted on $(T-\hat\lambda_k)Y$ there.
\end{ins}
\item p.~35, proof of Lemma~25: ``$\hat p_k$ is the mean of $n$ Bernoulli($p_0$)
  variables''.
\item p.~35, Lemma~26(v) and the sentence after it: see M24.
\end{itemize}
\note{The number is $\theta_0(1-p_0)/\{p_0(N-n)\}=3\cdot0.5/(0.5\cdot160)=0.019$ for
$N=200$ and $K=5$; with Chang's $K=2$ it is $0.030$. The SD of 0.025 is from your
Table~1.}

\paragraph{M5 (pp.~8--9). $\tilde\theta$ and $\hat\theta$.}
Definition~2(c) calls the final estimator $\hat\theta$, the limit theorem on
p.~9 uses $\hat\theta$, but the variance estimators, Propositions~7--8 and the
whole appendix use $\tilde\theta$ for it, and Section~5 switches back to
$\hat\theta$. Use $\tilde\theta$ everywhere, as Chang does ($\tilde\theta_k$ for
the folds, $\tilde\theta$ for their average).

\paragraph{M6 (pp.~8--9). Assumption 3 and the limit theorem.}
\begin{itemize}
\item $\Sigma_{10}$ has a misplaced bracket. It should be
  $\Sigma_{10}\equiv\E\big[\big(\psi_1(W,\theta_0,p_0,\eta_{10})+G_{1p_0}(D-p_0)\big)^2\big]$.
\item (a) ``requires that every subject has a bounded, non-zero probability of
  receiving a propensity score'' should be ``requires that, for every value of
  the covariates, the probabilities of being treated and of being untreated are
  both at least $\kappa$''.
\item (b) $\|UV_1\|_{P,4}\le C$ is a moment bound on the product of the two
  residuals, not a ``correlation''.
\item (e) ``The variance of the variance of the score'' should be ``The asymptotic
  variance must be positive''.
\item (f)(ii) ``rules out that the denominator in the score function to go
  crazy'' should be ``keeps the denominator $1-\hat g_k$ away from zero''.
\item (f)(iii): the sentence says the sum must vanish ``at a rate of
  $N^{-1/2}$'', but the display bounds it by $\varepsilon_N^2$. The rate
  $N^{-1/2}$ only follows once $\varepsilon_N=o(N^{-1/4})$. Write ``must be at
  most $\varepsilon_N^2$''.
\item The sentence on repeated cross-sections (``and the error between the
  error term'') is garbled. Write: ``For repeated cross-sections the conditions
  are analogous, with two additions: $\E[Y^2\mid X]\le C$ and $|\E[YU]|\le C$.
  The exact wording is in the appendix.''
\item The limit theorem should be attributed and use convergence in
  distribution:
\end{itemize}
\begin{ins}
\textcite[Theorems~3.1 and~3.2]{chang2020double} shows that if these conditions
hold with $\varepsilon_N=o(N^{-1/4})$, then
\[
  \sqrt N(\tilde\theta-\theta_0)\xrightarrow{d}N(0,\Sigma),
\]
with $\Sigma=\Sigma_{10}$ for repeated outcomes and $\Sigma=\Sigma_{20}$ for
repeated cross-sections, and that the variances are consistently estimated by
\begin{align*}
  \hat\Sigma_1&=\frac1K\sum_{k=1}^K\E_{n,k}\Big[\big(\psi_1(W,\tilde\theta,\hat p_k,\hat\eta_{1k})
     +\hat G_{1p}(D-\hat p_k)\big)^2\Big],\\
  \hat\Sigma_2&=\frac1K\sum_{k=1}^K\E_{n,k}\Big[\big(\psi_2(W,\tilde\theta,\hat p_k,\hat\lambda_k,\hat\eta_{2k})
     +\hat G_{2p}(D-\hat p_k)+\hat G_{2\lambda}(T-\hat\lambda_k)\big)^2\Big],
\end{align*}
where $\E_{n,k}[f(W)]=n^{-1}\sum_{i\in I_k}f(W_i)$,
$\hat G_{1p}=\hat G_{2p}=-\tilde\theta/\hat p_k$, and $\hat G_{2\lambda}$ is the
sample analogue of
\[
  G_{2\lambda_0}=\E[\partial_\lambda\psi_2]
  =-\frac{(1-2\lambda_0)\theta_0+\E\big[Y(D-g_0)/\{p_0(1-g_0)\}\big]}{\lambda_0(1-\lambda_0)} .
\]
\end{ins}
\note{$\hat G_{2\lambda}$ is used but never defined in v13. I derived the
expression above from $\psi_2$; it agrees with the display in your appendix
(p.~41). Check that your code uses the same expression.}

\paragraph{M7 (p.~10). Section 3.3.}
\begin{itemize}
\item Say that $q_i$ contains a constant whose coefficient is not penalised, in
  both regressions. Lemma~10 and the pooling result rely on it.
\item ``the penalty $\lambda$ can be either cross-validated or implemented in a
  plug-in form Belloni and Chernozhukov (2011)'': the plug-in rule for the logit
  Lasso is from \textcite{belloni2016post}, and the rule with loadings for the
  outcome Lasso is from \textcite{belloni2012sparse}, which is also what Chang
  cites.
\item The loadings are called $\gamma_{kj}$, but $\gamma$ is the vector of
  propensity coefficients from Section~4.1 on. Call the loadings
  $\hat\psi_{kj}$, as in Section~5.
\item Last sentence: the reference for the refit is
  \textcite{belloni2013least}, not Belloni et al.\ (2014).
\end{itemize}
Replacement for the last four sentences of Section~3.3:
\begin{ins}
Chang plugs the penalised fits into the score directly, without refitting; in
his application, the penalty is chosen by cross-validation. In this thesis, I
instead refit each nuisance without a penalty on its own selected support
(post-Lasso; \cite{belloni2013least}). The refit removes the shrinkage of the
selected coefficients, which turns out to be a large source of bias in
Section~\ref{sec:mc}, where I report both versions. Neither version is
post-double-selection \parencite{belloni2014inference}, the estimator that
\textcite{kolesar2026fragility} study. In post-double-selection, the covariates
selected by a Lasso of the outcome and by a Lasso of the treatment are pooled and
enter one OLS regression. In DML-DiD, each fitted nuisance enters the orthogonal
score as a prediction, and cross-fitting replaces the OLS step.
\end{ins}

\subsection{Section 4, up to Proposition 8 (pp.~10--14)}

\paragraph{M8 (p.~11).} The sentence ``For DML-DiD, it is enough \dots\ product of
the two sparsity indices'' appears twice. Delete the second one (``For DML-DiD,
it is enough to require \dots\ Smucler et al., 2019)'') and keep ``I will start by
investigating \dots''.

\paragraph{M9 (p.~12). Lemma 4.}
The first ``$\equiv$'' should be ``$=$'' (it is the claim, not a definition), a
closing bracket is missing, and the bound should have $\kappa'$, the constant of
the lemma, not $\kappa$. ``6 below translates'' should be ``Proposition~6 below
translates''.
\begin{ins}
\begin{lemma}[Mixed bias] % \label{lem:bias}
Let $g$ and $\ell$ be fixed functions with $\kappa'\le g\le1-\kappa'$ for some
$\kappa'\in(0,1/2]$. Then
\[
  \E[\psi_1(W,\theta_0,p_0,g,\ell)]=B_1(g,\ell)\equiv
  \E\!\left[\frac{(\ell-\ell_{10})(g-g_0)}{p_0(1-g)}\right],
  \qquad
  |B_1(g,\ell)|\le\frac{\|g-g_0\|\,\|\ell-\ell_{10}\|}{p_0\kappa'} .
\]
For repeated cross-sections, the same holds with $\psi_2$ and $\ell_{20}$ in
place of $\psi_1$ and $\ell_{10}$, and an additional factor
$1/\{\lambda_0(1-\lambda_0)\}$.
\end{lemma}
\end{ins}
\note{``The proof can be found in the appendix'': I could not find a proof of
Lemma~4 in the appendix of v13. Here is one; put it before Lemma~25.}
\begin{ins}
\begin{proof}[Proof of Lemma~\ref{lem:bias}]
Write $\Delta Y=Y(1)-Y(0)$ and $m_1(X)=\E[\Delta Y\mid X,D=1]$, the trend of the
treated, so that $\ell_{10}(X)=\E[\Delta Y\mid X,D=0]$ is the trend of the
untreated. The two terms of $\psi_1$ combine into
$\psi_1(W,\theta,p_0,g,\ell)=(\Delta Y-\ell)(D-g)/\{p_0(1-g)\}-\theta$, so
\begin{equation*}
  p_0\,\E[\psi_1(W,\theta_0,p_0,g,\ell)]
  =\underbrace{\E\Big[(\Delta Y-\ell)\frac{D-g}{1-g}\Big]}_{\text{(a)}}-\underbrace{p_0\theta_0}_{\text{(b)}} .
\end{equation*}
\emph{Term (a).} Conditional on $X$, a unit is treated with probability $g_0$.
For a treated unit, $(D-g)/(1-g)=1$ and $\E[\Delta Y-\ell\mid X,D=1]=m_1-\ell$.
For an untreated unit, $(D-g)/(1-g)=-g/(1-g)$ and
$\E[\Delta Y-\ell\mid X,D=0]=\ell_{10}-\ell$. Hence
\[
  \E\Big[(\Delta Y-\ell)\frac{D-g}{1-g}\,\Big|\,X\Big]
  =g_0(m_1-\ell)-(1-g_0)\frac{g}{1-g}(\ell_{10}-\ell).
\]
\emph{Term (b).} By RO(i) and RO(iii), $m_1-\ell_{10}=\E[Y^1(1)-Y^0(1)\mid X,D=1]$,
the treatment effect on the treated at $X$. Since
$\E[f(X)\mid D=1]=\E[Df(X)]/p_0=\E[g_0f(X)]/p_0$ for any $f$,
\[
  p_0\theta_0=p_0\,\E[m_1-\ell_{10}\mid D=1]=\E\big[g_0(m_1-\ell_{10})\big].
\]
This is the expectation of the display for term (a) at the true nuisances,
$(g,\ell)=(g_0,\ell_{10})$, where the second part vanishes; equivalently,
$\E[\psi_1(W,\theta_0,p_0,g_0,\ell_{10})]=0$.

\emph{Subtracting (b) from (a).} Both terms are now expectations of functions of
$X$. In the difference, the unknown trend of the treated cancels,
$g_0(m_1-\ell)-g_0(m_1-\ell_{10})=g_0(\ell_{10}-\ell)$, so
\[
  p_0\,\E[\psi_1(W,\theta_0,p_0,g,\ell)]
  =\E\Big[(\ell_{10}-\ell)\Big(g_0-\frac{(1-g_0)g}{1-g}\Big)\Big]
  =\E\Big[\frac{(\ell_{10}-\ell)(g_0-g)}{1-g}\Big],
\]
where the last step uses $g_0(1-g)-(1-g_0)g=g_0-g$. Since
$(\ell_{10}-\ell)(g_0-g)=(\ell-\ell_{10})(g-g_0)$, dividing by $p_0$ gives the
identity. The bound follows from $1-g\ge\kappa'$ and the Cauchy--Schwarz
inequality. For repeated cross-sections, replace $\Delta Y$ by
$(T-\lambda_0)Y/\{\lambda_0(1-\lambda_0)\}$ and $\ell$ by
$\ell/\{\lambda_0(1-\lambda_0)\}$.
\end{proof}
\end{ins}

\paragraph{M10 (p.~12). The paragraph after Assumption 5.}
The paragraph cites results for the plain Lasso with Gaussian errors
(Belloni and Chernozhukov 2011, Corollary~1) and for a differently tuned logit
Lasso (van de Geer 2008). Your first steps are post-Lasso with data-driven
loadings and a logit post-Lasso with the plug-in penalty, so cite the results
for those estimators. Also ``the lasso first steps with root in this basis'' in
Assumption~5 should be ``the Lasso first steps computed in this basis''.
\begin{ins}
For the outcome nuisance, the assumption follows from
\textcite[Theorems~1 and~2]{belloni2012sparse}, for the Lasso and for post-Lasso,
which allow non-Gaussian, heteroskedastic errors with the data-driven penalty
loadings used in Section~\ref{sec:mc}; with Gaussian, homoskedastic errors, see
\textcite[Corollaries~1 and~3]{belloni2011high}. These rates are stated in the
empirical prediction norm on the estimation sample. They carry over to the
population norm because the post-Lasso error is sparse, with at most $\hat s+s$
non-zero coordinates and $\hat s\lesssim s$, and because, with probability
approaching one, the sparse eigenvalues of the empirical Gram matrix are within
constant factors of the population ones, for Gaussian as well as bounded designs
\parencite[Lemmas~1 and~2]{belloni2011high}. The outcome Lasso uses only control
observations, so these conditions and the rate refer to the covariate
distribution of the controls. Under overlap, the rate carries over to the full
population, since $dP_X/dP_{X\mid D=0}=(1-p_0)/(1-g_0)\le(1-p_0)/\kappa$. For the
propensity score, the assumption follows from
\textcite[Appendix~B.2, Lemmas~1 and~2]{belloni2016post}, for the logit Lasso and
its post-Lasso refit. Their rate is in the prediction norm weighted by
$g_0(1-g_0)\ge\kappa(1-\kappa)$, and since $|\Lambda(a)-\Lambda(b)|\le|a-b|/4$, it
implies the rate for $\hat g_k$. Their result needs an identifiability
condition, which for general bounded designs holds if
$s_g^2\log(p\vee N)=o(N)$. The assumed rates are upper bounds on the estimation
error; smaller sparsity indices give faster rates.
\end{ins}
\note{I checked these theorem and lemma numbers against the arXiv versions
(1010.4345, 1106.5242, 1304.3969). Check them against the published versions
before you submit, since numbering sometimes changes.}

\paragraph{M11 (pp.~12--13). Proposition 6 and the text around it.}
\begin{itemize}
\item ``for $c=\{g,\ell\}$'' should be ``for $c\in\{g,\ell\}$'', and
  $\hat\ell_k$ should be $\hat\ell_{1k}$ as elsewhere.
\item ``Conversely, 6(b) follows by \dots'': it is not a converse. Write
  ``Part~(b) follows by multiplying the two rates.''
\item ``this implies $s_g\log[p]=s_\ell\log[p]=o(\sqrt N)$'' asserts that the two
  indices are equal. Write ``this requires $s_g\log p=o(\sqrt N)$ and
  $s_\ell\log p=o(\sqrt N)$''.
\item The product condition is not new and the text should say so, see
  Part~\ref{sec:obvious}. Add after the proposition:
\end{itemize}
\begin{ins}
The product condition is the one \textcite[Remark~5.2]{chernozhukov2018double}
give for the average treatment effect on the treated in a cross-section, and the
one \textcite{farrell2015robust} and \textcite{smucler2019unifying} obtain for
doubly robust estimators: if the propensity score is very sparse, the outcome
regression may be quite dense, and vice versa.
\textcite[footnote~13]{kolesar2026fragility} note that it replaces their
condition \eqref{eq:sparsity2} under cross-fitting, and that such assumptions on
the propensity score are likewise sensitive to the normalisation of the
controls. The rest of this section works out that remark for DML-DiD.
\end{ins}

\paragraph{M12 (p.~13). Proposition 7 (wrong as printed in (b)).}
In (b), ``$0\le1$'' has lost the $h$; the upper bound on $g_N$ is $1-\kappa/2$,
not $1/2-\kappa/2$; the two norms are deterministic, so $O$, not $O_P$; and the
proof shows $\{1+o_P(1)\}$, not $\{1-o_P(1)\}$. In (a), ``for all $k\in K$''
should be ``for all $k$''. In (c), $z_{1-a/2}$ should be $z_{1-\alpha/2}$.
\begin{ins}
\begin{proposition}[Rate failures] % \label{prop:ratefail}
Suppose RO(i)--RO(iii) and Assumption~\ref{ass:chang}(a)--(e) hold, and let
$b,c>0$ with $b+c<1/2$.
\begin{enumerate}
\item[(a)] If, with probability $1-o(1)$, $\|\hat g_k-1/2\|_{P,\infty}\le1/2-\kappa$
  for all $k$, and $\|\hat g_k-g_0\|=O_P(N^{-b})$ and
  $\|\hat\ell_{1k}-\ell_{10}\|=O_P(N^{-c})$, then
  $\tilde\theta-\theta_0=O_P(N^{-(b+c)})$.
\item[(b)] Let $h=h(X)$ satisfy $0\le h\le1$ and $\E[h^2]>0$, and let
  $g_N=g_0+N^{-b}h$ and $\ell_N=\ell_{10}+N^{-c}h$. Then
  $\|g_N-g_0\|=O(N^{-b})$, $\|\ell_N-\ell_{10}\|=O(N^{-c})$, and
  $\kappa/2\le g_N\le1-\kappa/2$ for $N\ge(2/\kappa)^{1/b}$. If
  $(\hat g_k,\hat\ell_{1k})=(g_N,\ell_N)$ for all $k$, then
  \[
    \tilde\theta-\theta_0=B_1(g_N,\ell_N)\{1+o_P(1)\},\qquad
    B_1(g_N,\ell_N)\ge N^{-(b+c)}\,\E[h^2]/p_0 .
  \]
\item[(c)] In the setting of (b),
  $\Pr\big(\theta_0\in\big[\tilde\theta\pm z_{1-\alpha/2}(\hat\Sigma_1/N)^{1/2}\big]\big)\to0$.
\end{enumerate}
\end{proposition}
The proof is in the appendix. Part~(a) says that the error of DML-DiD is at most
of the order of the product of the two first-step errors. Parts~(b) and~(c) show
that this order can be attained, and that the confidence interval then misses
$\theta_0$ with probability approaching one. The sequence in (b) is a worst
case: the two errors are proportional to the same function $h$, so their product
has one sign everywhere. Errors of an actual Lasso need not line up in this way.
Proposition~\ref{prop:stein} below shows what lining up means for coefficient
vectors, and D2 and D3 compare the two cases.
\end{ins}

\paragraph{M13 (p.~13). The paragraph after Proposition 8.}
``This proposition assesses what happens \dots\ when the nuisance estimators both
converge slowly'' (after Proposition~7) is imprecise: one of them may converge
fast; what matters is $b+c$. After Proposition~8, ``in the proof of'' starts a
sentence in lower case, ``27(ii)'' needs ``Lemma'', and what uses
$\hat\Sigma_1=O_P(1)$ is Proposition~7(c). Replacement:
\begin{ins}
The proof is Lemma~\ref{lem:variance} in the appendix.
Propositions~\ref{prop:ratefail} and~\ref{prop:ratedr} are two sides of the same
calculation. In terms of the rates $N^{-b}$ and $N^{-c}$ of the two first steps,
$b+c>1/2$ suffices for a valid limit distribution, while $b+c<1/2$ can make
coverage go to zero. At $b+c=1/2$ the bias is of the same order as the standard
error, so coverage is distorted but need not go to zero.

Proposition~\ref{prop:ratedr} concerns the estimator. For the interval
$\tilde\theta\pm z_{1-\alpha/2}(\hat\Sigma_1/N)^{1/2}$ to have asymptotic coverage
$1-\alpha$, the variance estimator must also be consistent. In the proof of
\textcite[Theorem~3.2]{chernozhukov2018double}, the error of the variance
estimator depends only on how fast each first step converges, not on the product
of their errors, so consistency of $\hat g_k$ and $\hat\ell_{1k}$ is enough. I do
not prove this here. Lemma~\ref{lem:consist}(ii) shows the weaker statement
$\hat\Sigma_1=O_P(1)$, which is what Proposition~\ref{prop:ratefail}(c) needs.

Proposition~\ref{prop:ratedr} is the result of
\textcite[Theorem~5.1]{chernozhukov2018double} for the average treatment effect
on the treated, applied to the change in outcomes: $\psi_1$ is their score (5.4)
with $Y$ replaced by $Y(1)-Y(0)$, except that $\theta$ enters as $-\theta$
instead of $-D\theta/p$. Their rate condition for this score involves only the
product of the two errors. Chang instead uses their general condition (3.7)--(3.8),
which bounds every error by one sequence $\varepsilon_N=o(N^{-1/4})$ and, in his
Assumption~3.1, adds the squared propensity error. Lemma~\ref{lem:bias} shows
that the squared propensity error does not enter the bias of $\psi_1$, so the
weaker condition suffices. The product condition guarantees validity in the
limit; D2 shows that close to the boundary $b+c=1/2$ convergence can be too slow
to be visible at realistic sample sizes.
\end{ins}
\note{This replaces the two paragraphs from ``The proof is in the appendix.
Proposition 7 and proposition 8 are two sides \dots'' to ``\dots bounds all rates
by a single $\varepsilon_N$'' (pp.~13--14).}

\subsection{Section 4.1: non-sparse data (pp.~14--15)}

\paragraph{M14 (p.~14). The proof of the Remark has a wrong step (wrong as printed).}
``the numerator simplifies to
$\mathrm{Cov}(g_0,\ell_{10})=p_0/(1-p_0)(\E[\ell_{10}\mid D=0])$'' is not true.
Also make the Remark a numbered environment: ``as per~.'' on p.~20 is an empty
reference because the Remark has no number.
\begin{ins}
\begin{proof}
By Lemma~\ref{lem:bias},
$B_1(p_0,c)=\E[(c-\ell_{10})(p_0-g_0)]/\{p_0(1-p_0)\}$. Since $\E[g_0]=p_0$, the
constant $c$ drops out and the numerator equals
$\operatorname{Cov}(g_0,\ell_{10})$. Further,
$\E[g_0\ell_{10}]=\E[D\ell_{10}]=p_0\E[\ell_{10}\mid D=1]$ and
$\E[\ell_{10}]=p_0\E[\ell_{10}\mid D=1]+(1-p_0)\E[\ell_{10}\mid D=0]$, so
$\operatorname{Cov}(g_0,\ell_{10})=p_0(1-p_0)\{\E[\ell_{10}\mid D=1]-\E[\ell_{10}\mid D=0]\}$.
The last claim is Stein's lemma, $\E[Xf(X)]=\E[\nabla f(X)]$ for
$X\sim N(0,I_p)$, applied to $f(X)=\Lambda(a+X'\gamma)$.
\end{proof}
\end{ins}

\paragraph{M15 (pp.~14--15). Proposition 9: notation, and a simpler form.}
\begin{itemize}
\item $w$ is used both for the function $w(X)$ and for its mean, but
  $\E[Xw]$ and $\nabla w$ need the function.
\item $c_1$ and $c_2$ are defined as expectations, but the proof writes
  ``$=c_1\gamma+c_2\Delta\gamma$'' for the expression inside the expectation.
\item The third line of the proof has $\Lambda(a-X'\gamma)$; it is
  $\Lambda(a+X'\gamma)$.
\item $\gamma\Delta\delta$ and $\Delta\gamma\Delta\delta$ are inner products and
  need transposes.
\item $c_1,c_2$ are also the names of the adjustment terms in $\psi_1,\psi_2$
  (p.~7) and of the three terms in Proposition~12.
\end{itemize}
While checking the proof I noticed that the two constants simplify, because
$\Lambda'=\Lambda(1-\Lambda)$:
\[
  \Lambda'(\hat z)(1-g_0)-\Lambda'(z)(1-\hat g)=(1-\hat g)(1-g_0)(\hat g-g_0),
  \qquad
  \frac{\Lambda'(\hat z)(1-g_0)}{(1-\hat g)^2}=\frac{\hat g(1-g_0)}{1-\hat g}.
\]
With this, the constants can be read: $\bar c_2=1$ and $\bar c_1=0$ when
$\hat g=g_0$, and the statement about ``$c_1$ is small'' becomes exact. I
recommend this version (statement, proof and the text after it, replacing
everything from Proposition~9 to ``\dots can cancel in case not.''):
\begin{ins}
\begin{proposition}[Bias in the Gaussian design] % \label{prop:stein}
In this design,
\[
  B_1(\hat g,\hat\ell)=\bar c_2\,\Delta\gamma'\Delta\delta+\bar c_1\,\gamma'\Delta\delta+(\hat b-b)\,\bar w ,
\]
where, with $w(X)=(\hat g-g_0)/\{p_0(1-\hat g)\}$,
\[
  \bar w=\E[w(X)],\qquad
  \bar c_1=\E\big[(1-g_0)\,w(X)\big],\qquad
  \bar c_2=\E\left[\frac{\hat g\,(1-g_0)}{p_0(1-\hat g)}\right].
\]
\end{proposition}
\begin{proof}
By Lemma~\ref{lem:bias} and $\hat\ell-\ell_{10}=(\hat b-b)+X'\Delta\delta$,
$B_1=\E[(\hat\ell-\ell_{10})\,w(X)]=(\hat b-b)\bar w+\Delta\delta'\E[X\,w(X)]$, and
by Stein's lemma $\E[X\,w(X)]=\E[\nabla w(X)]$. Since
$\hat g-g_0=(1-g_0)-(1-\hat g)$, we can write
$w=p_0^{-1}\{(1-g_0)/(1-\hat g)-1\}$. With
$\nabla\hat g=\hat g(1-\hat g)\hat\gamma$ and $\nabla g_0=g_0(1-g_0)\gamma$,
\[
  \nabla w=\frac{(1-g_0)\hat g(1-\hat g)\hat\gamma-g_0(1-g_0)(1-\hat g)\gamma}{p_0(1-\hat g)^2}
  =\frac{(1-g_0)\,(\hat g\,\hat\gamma-g_0\,\gamma)}{p_0(1-\hat g)} .
\]
Substituting $\hat\gamma=\gamma+\Delta\gamma$ gives
$\nabla w=\frac{(1-g_0)(\hat g-g_0)}{p_0(1-\hat g)}\gamma+\frac{\hat g(1-g_0)}{p_0(1-\hat g)}\Delta\gamma$,
and taking expectations gives the result.
\end{proof}
The three terms have a clear reading. The first is the inner product of the two
coefficient errors. The second and third are weighted averages of the signed
propensity error $\hat g-g_0$; they vanish when the propensity score is correct,
where $\bar c_2=\E[g_0]/p_0=1$. If both first steps fit only intercepts, then
$\Delta\gamma=-\gamma$, $\Delta\delta=-\delta$, $\bar w=0$, $\bar c_2=1$ and
$\bar c_1=\operatorname{Var}(g_0)/\{p_0(1-p_0)\}$, and the proposition
reproduces the Remark. When $\hat g$ is close to $g_0$, the bias is therefore
approximately the inner product of the two coefficient errors,
\begin{equation}\label{eq:inner}
  B_1\approx\bar c_2\sum_{j=1}^p\Delta\gamma_j\,\Delta\delta_j .
\end{equation}
For post-Lasso, the first step used in this thesis, the refit estimates the
selected coefficients without shrinkage, so
$\Delta\gamma_j\approx-\gamma_j\mathbf 1\{j\notin\hat S_g\}$,
$\Delta\delta_j\approx-\delta_j\mathbf 1\{j\notin\hat S_\ell\}$ and
\begin{equation}\label{eq:both}
  B_1\approx\bar c_2\sum_{j\notin\hat S_g\cup\hat S_\ell}\gamma_j\,\delta_j .
\end{equation}
Only covariates that \emph{both} fits omit contribute. This is the DML-DiD
analogue of the double under-selection that \textcite{wuthrich2023omitted}
identify for post-double-selection. Without the refit, as in Chang's
specification, products of the shrinkages of the selected coefficients are
added. In DiD terms, $\gamma_j$ measures how covariate $j$ drives selection into
treatment and $\delta_j$ how it drives the untreated trend. The bias comes from
covariates that do both and that both fits miss. It is large when the products
$\gamma_j\delta_j$ share a sign (\emph{aligned} errors, as in
Proposition~\ref{prop:ratefail}(b)) and largely cancels otherwise. D2 and D3
test \eqref{eq:inner} and \eqref{eq:both}.
\end{ins}
\note{Your code computes $c_1$ and $c_2$ in the old form. The two forms are
algebraically identical, so the numbers in the tables do not change.}

\subsection{Section 4.2: the basis (pp.~15--19)}

\paragraph{M16 (p.~16). Lemma 10.}
The statement ends ``are in the same bases, $q$ and $\tilde q$'' and has
$\hat g_K$; ``hold the number of folds fixed'' should be ``hold the folds
fixed''; the first line of the proof stops in the middle; the sentence about the
Lasso belongs after the proof; and the lemma needs the intercept and the
existence of the logit estimate. The caveats are needed in Section~6, where you
drop a product group that has no controls.
\begin{ins}
\begin{lemma}[Invariance of unpenalised first steps] % \label{lem:invariance}
Let $q_i$ contain an intercept, let $\tilde q=Aq$ with $A$ invertible, and hold
the folds fixed. Let $\hat g_k$ be the unpenalised logit fit and $\hat\ell_{1k}$
the OLS fit on the controls, both computed on $I_k^c$, and suppose the logit
estimate exists. Then $\hat g_k$ and $\hat\ell_{1k}$, and hence $\tilde\theta$,
are the same in the bases $q$ and $\tilde q$.
\end{lemma}
\begin{proof}
For every $\beta$, $q'\beta=\tilde q'\tilde\beta$ with
$\tilde\beta=(A')^{-1}\beta$. The logit log-likelihood and the least-squares
criterion depend on the coefficients only through the index $q'\beta$. Hence
$\tilde\beta$ minimises the criterion in the basis $\tilde q$ exactly when
$A'\tilde\beta$ minimises it in the basis $q$, and the fitted values coincide.
\end{proof}
The Lasso adds the penalty $\lambda\sum_j\hat\psi_j|\tilde\beta_j|$, which is not
a function of the index alone, so the argument fails. The same argument shows
that a Lasso fit is invariant to a change of basis that only involves columns
whose coefficients are not penalised, for example a change of reference
category among dummies that are forced into the model. Lemma~\ref{lem:invariance}
is the DiD counterpart of the observation of \textcite{kolesar2026fragility} that
OLS is invariant to the representation: with unpenalised first steps,
$\mathcal R(\mathcal A)=0$. There are two caveats. The logit estimate does not
exist when some category contains only treated or only untreated units, and when
the number of columns is not small relative to the sample it is badly biased
\parencite{sur2019modern}. Outcome-regression DiD (OR-DiD), which needs only OLS
among the controls, is invariant whenever the controls outnumber the columns of
the dictionary and every category contains controls.
\end{ins}

\paragraph{M17 (p.~16). Lemma 11.} Correct. Two notational points: the
difference of the two fits is called $\delta_g$ in the statement and $\delta$ in
the bound and in the sentence after it, and $\delta$ is the vector of trend
coefficients everywhere else. Call it $\Delta g$, as Proposition~12 does. ``Let
$g,g'<1$'' can simply be ``Let $g,g'\in[\kappa,1-\kappa]$'' from the start.

\paragraph{M18 (pp.~16--17). Proposition 12 (wrong as printed: ``$=$'' for ``$\le$'').}
\begin{itemize}
\item The bound is written with ``$=$''. It is an inequality, and the sum needs
  brackets around all three terms:
  \begin{multline*}
  \big|\tilde\theta^{(A)}-\tilde\theta^{(A_0)}\big|\le\max_k\frac1{\hat p_k}\frac1n\sum_{i\in I_k}
  \bigg(\|\Delta g_k\|_{P,\infty}\frac{(1-D_i)\big|\Delta Y_i-\hat\ell_k^{(A_0)}\big|}{\kappa^2}\\
  +\|\Delta g_k\|_{P,\infty}\frac{(1-D_i)|\Delta\ell_k|}{\kappa^2}+\frac{|\Delta\ell_k|}{\kappa}\bigg).
  \end{multline*}
\item The title ``Estimator bias in response to basis shift'' and ``the
  deviance in the estimators'' and ``The bound of the bias'': nothing here is a
  bias. It is the \emph{difference between two estimates}. Title: ``Change of the
  estimate under a change of basis''.
\item ``$\Delta\ell_k\equiv\hat\ell_j^{(A)}-\hat\ell_j^{A_0}$, $j=\{1,2\}$'' should
  be $\Delta\ell_k\equiv\hat\ell_{1k}^{(A)}-\hat\ell_{1k}^{(A_0)}$.
\item The three terms are named $c_1,c_2,c_3$; see M15.
\item In $\mathcal R(\mathcal A)$, say which variance: ``$\hat\Sigma_1^{(A_0)}$,
  the variance estimate in the reference basis $A_0$''.
\item ``If the rate conditions hold uniformly over a fixed and finite set of
  bases, the first order terms in (12) vanish, and
  $\mathcal R(\mathcal A)=o_p(1)$'': the conclusion is right, but it does not
  follow from Proposition~12. It follows from Lemma~26. Replacement:
\end{itemize}
\begin{ins}
If the conditions of Proposition~\ref{prop:ratedr} hold in every basis of a
fixed, finite set $\mathcal A$, then by Lemma~\ref{lem:variance} every
$\tilde\theta^{(A)}$ has the same expansion
$N^{-1}\sum_i\varphi(W_i)+o_P(N^{-1/2})$, which does not depend on the basis.
The estimates then differ by $o_P(N^{-1/2})$, and
$\mathcal R(\mathcal A)=o_P(1)$. A range that is not small relative to one
standard error is therefore evidence that the rate conditions, and in this
context sparsity, fail in at least one of the bases.
\end{ins}

\paragraph{M19 (p.~17). ``INDSÆT HER''.} The running example goes here; the text
and the table are in Part~\ref{sec:pool-ins}. ``To formalise this insight,
consider the following lemma'' introduces a proposition.

\paragraph{M20 (p.~17--18). Proposition 13 (wrong as printed).}
\begin{itemize}
\item The sum over the pools is missing in the display. As printed, $B_1$ is
  the contribution of a single, unspecified pool $G$.
\item ``suppose both full and threshold coding, and that'' is not a sentence.
  It should be ``suppose, as holds for reference and threshold coding, that''.
\item ``if their rows are similar'' should be ``if their rows agree''.
\item The proof has a sign error:
  $\E[(\hat\ell-\ell_{10})(\hat g_G-g_0)\mid Z\in G]$ equals
  $\operatorname{Cov}(\ell_{10}-\hat\ell,g_0\mid Z\in G)$, not
  $\operatorname{Cov}(\hat\ell-\ell_{10},g_0\mid Z\in G)$. (The statement has the
  right sign.)
\item ``with a reference category spans all functions of $Z$'': it is the
  intercept together with the columns that spans them.
\end{itemize}
The corrected statement and proof are in Part~\ref{sec:pool-ins}.

\paragraph{M21 (p.~18). Corollary 14 and the two remarks.}
The corollary and its proof are correct (``If it atypical in both'' needs
``is''). Remark~1 ends with a wrong conclusion \textbf{(wrong as printed)}:
``An unlucky choice of category therefore matters when it is large, i.e.\ when
it comprises a majority of the observations.'' The condition you derive in the
sentence before is that $\pi_b\sqrt N$ is not small. In D4a, a reference
category with 5\% of the sample gives coverage 0.77 at $N=4000$, and one with
15\% gives 0.06. Corrected remarks are in Part~\ref{sec:pool-ins}.

\subsection{Appendix (pp.~28--44)}

\paragraph{M22 (pp.~28--29). Chebyshev.}
Definition~18 has $t\sigma_X^2$; it is $t\sigma_X$. In the proof of Lemma~19,
``Since both $\Pr(|X-\E(X)|\ge t)$ and $\mathrm{Var}(X)/t^2$ are positive'' is not
the reason the two events are equal. The reason is that $t>0$ and
$x\mapsto x^2$ is increasing on $[0,\infty)$. (I recommend cutting this
material altogether, see Part~\ref{sec:obvious}.)

\paragraph{M23 (pp.~33--35). Lemma 25.} The decomposition and its proof are
correct. Fix: ``Lemma (4)'' twice (it prints like an equation number);
``$\hat p_k$ is the mean of $N-n$ Bernoulli($p_0$) variables'' should be $n$
(M4); and before ``On the event that $\{\hat p_k\ge p_0/2\}$'' add ``whose
probability tends to one, since $\hat p_k\xrightarrow{p}p_0>0$''. The step
$\E[\psi_1^2]\le\E[(\psi_1+\theta_0)^2]$ is true only because $\E[\psi_1]=0$;
say so.

\paragraph{M24 (p.~35). Lemma 26.}
\begin{itemize}
\item The title is ``The variance estimator''; the lemma is asymptotic
  linearity.
\item ``Fix number of folds $k$'' should be ``Fix a fold $k$''. In (i) a bracket
  is missing. In (ii) ``On the event that
  $\Pr(\{\kappa'\le\hat g_k\le1-\kappa'\})\to1$'' should be ``On the event
  $\{\kappa'\le\hat g_k\le1-\kappa'\}$, whose probability tends to one''.
\item Step (v) is garbled, and the sentence after the display argues with
  auxiliary samples, which you no longer use for $\hat p_k$.
\end{itemize}
Replacement for (v) and the rest of the proof:
\begin{ins}
(v) Since
$p_0/\hat p_k-1=-(\hat p_k-p_0)/p_0+(\hat p_k-p_0)^2/(p_0\hat p_k)
=-(\hat p_k-p_0)/p_0+O_P(N^{-1})$ and $\hat p_k-p_0=\E_{n,k}[D-p_0]$,
\[
  \tilde\theta_k-\theta_0=\E_{n,k}\Big[\psi_1(W,\theta_0,p_0,\eta_{10})-\frac{\theta_0}{p_0}(D-p_0)\Big]+o_P(N^{-1/2})
  =\E_{n,k}[\varphi(W)]+o_P(N^{-1/2}).
\]
The folds partition the sample and have the same size $n=N/K$, so averaging over
$k$ gives $\tilde\theta-\theta_0=N^{-1}\sum_{i=1}^N\varphi(W_i)+o_P(N^{-1/2})$.
$\varphi$ has mean zero, since $\E[\psi_1(W,\theta_0,p_0,\eta_{10})]=0$ and
$\E[D]=p_0$. Its variance is finite by overlap and
Assumption~\ref{ass:chang}(d), and positive by
Assumption~\ref{ass:chang}(e). The central limit theorem completes the proof.
\end{ins}

\paragraph{M25 (pp.~36--37). Lemma 27.} The proof is correct. Fix:
\begin{itemize}
\item Statement: ``the set $\Pr(A_k(\kappa'))\to1$'' should be
  ``$\Pr(A_k(\kappa'))\to1$'', and $\delta_{g,k}\equiv\hat g_k-g_0$,
  $\delta_{\ell,k}\equiv\hat\ell_{1k}-\ell_{10}$ must be defined in the statement.
  There is a stray ``. :'' before (i).
\item First display of the proof: both terms should be multiplied by $p_0$, i.e.\
  ``$-p_0\theta_0$'' twice (they cancel, so nothing else changes).
\item ``by using lemma ??'' is Lemma~11.
\item In the display for $\E[(\Delta\psi)^2]$, the factor $1/p_0$ belongs to both
  terms and the second term has a minus sign:
  $\Delta\psi=p_0^{-1}\{(w_g-w_0)V_1-w_g\delta_\ell\}$. The bound in the next step
  is unaffected. The line with the absolute value around the whole expression is
  redundant.
\item ``On the set $A_k(\kappa')$, which per construction depends on the
  auxiliary sample $I_k^c$, the On $A_k(\kappa')$, which depends only on
  $I_k^c$'': delete the first half. ``since we assumed that
  $\|\delta_g\|,\|\delta_g\|=o_P(1)$'': the second is $\delta_\ell$.
\item (ii): ``Using the variance form in lemma 26'' should be ``Using the
  definition of $\hat\Sigma_1$''. ``$4/(\hat p_k^2\kappa'^2)$ and $2\tilde\theta$''
  should be $2\tilde\theta^2$. The citation ``(vandervaart)'' for ``it is enough
  to show that the variance within each fold is bounded'' is not needed: a
  finite sum of $O_P(1)$ terms is $O_P(1)$.
\end{itemize}

\paragraph{M26 (pp.~38--41). Proof of Proposition 7.} Correct. Fix:
\begin{itemize}
\item Part (a) starts ``By assumption, $P(A_k(\kappa'))\to1$''. The assumption of
  (a) is $A_k(\kappa)$.
\item ``in order to transfer it to condition 3.1(d) in Chang'': the band is
  part of condition (f) in Chang's Assumption~3.1, i.e.\ Assumption~3(f)(ii) in
  your numbering; (d) is the bound on $\E[V_1^2\mid X]$.
\item ``Where I have used that $0\le1-g_N\le1$'': you divide by $1-g_N$, so
  $0<1-g_N\le1$.
\item ``$B_1(g_N,\ell_N)=O(N^{-(b+c)})=o(1)$'': you have shown more, namely that
  $B_1(g_N,\ell_N)$ is between $N^{-(b+c)}\E[h^2]/p_0$ and
  $2N^{-(b+c)}\E[h^2]/(p_0\kappa)$. Say that; it is what (b) claims.
\item The display for the average over folds lacks brackets and the term
  $(p_0/\hat p_k-1)\theta_0$; write
  ``$\tilde\theta-\theta_0=B_1(g_N,\ell_N)+O_P(N^{-1/2})$ by averaging the
  previous display over $k$''.
\item Part (c): $T_N$ is first defined as an inequality (``$T_N=\sqrt N|\tilde\theta-\theta_0|\le\hat\sigma z$'').
  Define the statistic $T_N\equiv\sqrt N|\tilde\theta-\theta_0|$ and the coverage
  event $\{T_N\le z\hat\sigma\}$ separately.
\item ``The set $T_N\le M\cap\{|(N^{-(b+c)}\E[h^2/p_0])(1-\delta)|\}$ is only
  non-empty when \dots'' is not a set. Write: ``On
  $\{|\rho/B_1(g_N,\ell_N)|\le\delta\}$ we have
  $T_N\ge\sqrt N N^{-(b+c)}\E[h^2/p_0](1-\delta)$, so $\{T_N\le M\}$ can only
  occur if this lower bound is at most $M$ or if
  $|\rho/B_1(g_N,\ell_N)|>\delta$.''
\item $\delta$ is again a new object here; use $\epsilon'$ or $\tau$.
\end{itemize}

\paragraph{M27 (p.~41). The paragraph on repeated cross-sections stops without a conclusion.}
\begin{ins}
The argument carries over to repeated cross-sections with $\psi_2$, $\ell_{20}$
and $V_2$, and with the version of Lemma~\ref{lem:bias} that has the factor
$1/\{\lambda_0(1-\lambda_0)\}$. One term is new. Unlike $\hat p_k$, which
multiplies the whole fold average, $\hat\lambda_k$ enters the score itself, and
the score is not orthogonal in $\lambda$. Conditional on $I_k^c$, $\hat\lambda_k$
is fixed and shifts the conditional mean of the score by
$G_{2\lambda_0}(\hat\lambda_k-\lambda_0)+O_P(N^{-1})=O_P(N^{-1/2})$, because
$\hat\lambda_k$ is the mean of $N-n$ Bernoulli($\lambda_0$) variables. It also
moves the target of the outcome Lasso, $\E[(T-\hat\lambda_k)Y\mid X,D=0]$, away
from $\ell_{20}$ by $O_P(N^{-1/2})$. Since $b+c<1/2$, both are dominated by the
bias of order $N^{-(b+c)}$, and parts (a)--(c) hold as stated. Under the
conditions of Proposition~\ref{prop:ratedr}, the first of the two terms is of
the same order as the leading term; it is the reason for the term
$\hat G_{2\lambda}(T-\hat\lambda_k)$ in $\hat\Sigma_2$.
\end{ins}

\paragraph{M28 (pp.~42--44). Remaining.} ``On the $[\kappa,1.\kappa]$ bound''
should be $[\kappa,1-\kappa]$; ``using lemma ??'' is Lemma~11; in the regularity
conditions for repeated cross-sections, (a) should be
$\Pr(\kappa\le g_0(X)\le1-\kappa)=1$, ``$D=g_o(X)+U$'' should be $g_0$, and the
wording problems of M6 are repeated in (a), (b), (g) and (h).

% =============================================================================
\section{Inserts: title, introduction and Section 2}\label{sec:front}

\subsection{Title and section titles}
The title still reads ``My advisor told me sparsi ty was an assumption and he was
right'', with a broken word; Section~5 is ``Monte Carlo shows I'm right'' and
Section~6 ``Replication of a replication empirical application (say it fast
twice)''. Unless your supervisor has agreed to these, replace them. Suggestions:
\begin{itemize}
\item Title: \emph{Sparsity is a choice: representation and bias in double machine
  learning for difference-in-differences}. Plainer:
  \emph{Double machine learning for difference-in-differences when sparsity
  fails}.
\item Section~4: \emph{Sparsity, representation and bias} (now ``Enter
  sparsity'').
\item Section~5: \emph{Simulation evidence}.
\item Section~6: \emph{Chang's application revisited}.
\item Section~7: \emph{Conclusion}.
\end{itemize}

\subsection{Introduction (pp.~1--3)}

\where{p.~1, first sentence.}
The comparison of the two London water companies before and after one of them
moved its intake, which is the usual reference for the first
difference-in-differences, is in the second edition of Snow's book, from 1855.
The 1849 edition is a short pamphlet without it. Write ``Ever since John Snow's
study of cholera in London \parencite{snow1855mode}, \dots''.

\where{p.~1, second paragraph.}
``which is why several add-ons exist'' is informal (``several extensions
exist''), and ``heterogeneous treatment intensity, depending on the
characteristics of the treated unit'' does not describe
\textcite{callaway2024continuous}, which is about a treatment that varies in
intensity (dose) across units. Write ``or allow the treatment to vary in
intensity across units''.

\where{p.~2: ``He shows that the resulting estimator, the DML-DiD, is valid
whenever the first-step estimators converge faster than $N^{-1/4}$.''}
\begin{ins}
He shows that the resulting estimator, DML-DiD, is $\sqrt N$-consistent and
asymptotically normal when both first-step estimators converge faster than
$N^{-1/4}$.
\end{ins}

\where{p.~2: ``\textbf{kolesar<empty citation>} show that sparsity is \dots'' and
the next paragraph.}
The citation key is broken here (and for Abadie on pp.~6--7, Chang, Sequeira and
Kolesár on pp.~24--25, Chernozhukov on p.~37, Pitman on pp.~28--29, van der
Vaart on pp.~29 and~37). Search the source for \verb|{kolesar}|,
\verb|{abadie}|, \verb|{chang}|, \verb|{sequeira}|, \verb|{chernozhukov}|,
\verb|{pitman}| and \verb|{vandervaart}| and insert the full keys. In the same
paragraph: ``sparsity-based ML method'' needs a plural, ``choices that does not
affect'' should be ``do not'', and ``where the researcher choice enter'' should be
``where the researcher's choices enter''.

\where{p.~3, top: ``a representation that makes only one nuisances can be
harmless \dots\ in the untreated change in outcomes.''}
A word is missing in the second feature and the third is hard to follow.
\begin{ins}
Second, the bias depends on the \emph{product} of the two first-step errors: a
representation that makes only one nuisance non-sparse can be harmless, while
one that makes both non-sparse in the same direction is not. Third, the outcome
nuisance is a trend rather than a level, so what has to be sparse is how the
covariates shift the untreated \emph{change} in outcomes.
\end{ins}

\where{p.~3: the paragraph ``The thesis makes three contributions \dots'' and the
outline after it.}
Contribution (a) claims as new what is known (Part~\ref{sec:obvious}),
contribution (b) leaves out the pooling result, and the outline is off by one
section and promises a conclusion in ``section 6''.
\begin{ins}
The thesis makes three contributions. (a) It states the conditions for DML-DiD
in terms of the product of the two first-step errors. That the product is what
matters is known for doubly robust estimators
\parencite{chernozhukov2018double,smucler2019unifying}; I show that it replaces
Chang's requirement on each first step, and that coverage can go to zero when
the product condition fails. (b) It derives the bias that a failure of sparsity
creates. In a Gaussian design, the bias is the inner product of the two
first-step coefficient errors. With categorical covariates and post-Lasso first
steps, it is the covariance of selection and trends within the groups of
categories that the Lasso pools. This makes precise when the choice of reference
category matters: when the reference category differs from the categories it is
pooled with in \emph{both} selection and trends, in proportion to its share of
the sample. (c) It documents these mechanisms by simulation, and revisits
Chang's empirical application to measure how much choices of representation move
the estimate.

The rest of the thesis is organised as follows. Section~2 places the thesis in
the literature. Section~3 reviews the estimators of Abadie and Chang.
Section~4 contains the theory: what happens to DML-DiD when sparsity fails, and
how this depends on the representation of the covariates.
Section~\ref{sec:mc} provides simulation evidence. Section~\ref{sec:emp} revisits
Chang's application. Section~\ref{sec:conclusion} concludes.
\end{ins}
\note{If Section~6 ends up short or is cut, change (c) accordingly.}

\subsection{Section 2 (pp.~3--5)}

\where{p.~3: the linear model.}
The model $Y_i=D_i\theta+X_i'\beta+\varepsilon_i$, $D_i=X_i'\gamma+v_i$ is the
setting of post-double-selection, and of Kolesár et al. Say so, since the reader
otherwise wonders why a section on DiD starts with a cross-section:
``\dots consider the partially linear model that
\textcite{belloni2014inference} and \textcite{kolesar2026fragility} study''.
``Suppose that as a researcher you are interested in modeling the outcome
$Y_i$'' should be ``interested in $\theta$''.

\where{p.~3--4: ``Goodman-Bacon (2021) and Caetano and Callaway (2026) document
naïve linear inclusion \dots\ as e.g.\ formulated by Angrist (n.d.)''.}
A word is missing (``document \emph{that} naïve linear inclusion''), the book is
\textcite{angrist2009mostly}, and ``diverges from non-staggered treatment timing
with homogenous effects'' packs two different problems into one clause:
Goodman-Bacon is about timing, Caetano and Callaway about covariates. Split the
sentence in two.

\where{p.~4: ``Chang (2020) shows that plugging machine learning estimates \dots\
destroys $\sqrt N$-consistency. He restores it by deriving new score functions
of the outcome- and propensity regression \dots''}
\begin{ins}
\textcite{chang2020double} shows that plugging machine learning estimates of the
propensity score into the estimator of \textcite{abadie2005semiparametric}
does not in general give a $\sqrt N$-consistent estimator. He restores
$\sqrt N$-consistency by deriving Neyman-orthogonal scores, which add an outcome
regression to Abadie's moment condition, and by combining them with
cross-fitting.
\end{ins}

\where{p.~4: ``Ordinary Least Squares becomes infeasible, while sparsity-based
estimators remain valid.''}
Too strong: their validity is what the thesis questions. Write ``\dots OLS
becomes infeasible or very imprecise, while sparsity-based estimators remain
feasible''.

\where{p.~5: ``Smucler et al.\ (2019) show that \dots\ the final estimator will be
$\sqrt n$-consistent.''}
``two-nuisance influence function'' is their ``bilinear influence function'',
and the last sentence says the estimator is consistent ``as long as the other
nuisance [is] estimated consistently'', which is not enough: the other one must
be sparse enough to compensate.
\begin{ins}
\textcite{smucler2019unifying} show that for parameters with a bilinear
influence function, which include the ATT, cross-fitted estimators with
$\ell_1$-regularised first steps are $\sqrt n$-consistent and asymptotically
normal whenever the product of the sparsity indices of the two nuisances is
small enough, instead of each index on its own. The authors call this feature
rate double robustness: one nuisance may be far from sparse, provided it is
still estimated consistently and the other is sparse enough to compensate.
\end{ins}

\where{p.~5: Wüthrich and Zhu.} ``under half the chosen penalty'' is right. Add
what makes it relevant here: ``\dots can lead to bias in post-double-selection.
Proposition~\ref{prop:stein} gives the corresponding result for DML-DiD.''

% =============================================================================
\section{Inserts: categorical covariates and pooling (Section 4.2)}\label{sec:pool-ins}

\where{p.~17, from ``Proposition 12 shows how far the estimator moves \dots'' to
the end of Section~4 (p.~19). This replaces ``INDSÆT HER'', Proposition~13, its
proof, Corollary~14 and the two remarks.}
\note{I recomputed every number in the running example and in
Table~\ref{tab:pooling} today from \texttt{Bribes\_Regression.dta} in the
repository, with the sample definition of your notebook ($N=1084$). Recompute
them once in your own notebook (one \texttt{groupby} each), so that you can
answer questions about them. The product group variable is
\texttt{hc\_group}; transportation is group 14, machinery 13, animal products 1
and metals 12.}
\begin{ins}
\subsubsection*{Categorical covariates: what the Lasso pools}
Proposition~\ref{prop:decomp} describes how the estimate moves between bases.
With categorical covariates it is possible to say \emph{which} bases create
bias. As a running example, take the product group of a shipment in the data of
\textcite{sequeira2016corruption} that Chang uses. The shipments fall into 14
product groups. The largest, transportation (vehicles), accounts for 68.5\% of
the 1{,}084 shipments in Chang's sample, and it differs from the other groups in
both selection and trends: 5.1\% of its shipments are treated, against 40.4\% in
the other groups, and among untreated shipments the mean of
$\log(1+\text{bribe})$ is flat for transportation (0.91 before 2008, 0.88
after), while it falls from 7.95 to 1.01 in the other groups. The comparison
rests on 60 and 24 untreated shipments before 2008 and is only descriptive. It
shows, however, that the way product group enters the dictionary can matter.

A post-Lasso fit on dummies gives each category with a selected dummy its own
level, like a fixed effect, and gives all remaining categories one common level,
like pooled OLS. I say that these categories are \emph{pooled}. Just as pooled
OLS is biased when the ignored group effects are correlated with the regressor,
pooling biases DML-DiD when, among the pooled categories, those more often
treated also have different untreated trends. Table~\ref{tab:pooling}
illustrates how the reference category decides which categories are pooled.

\end{ins}
\begin{table}[ht]
\centering\small
\setlength{\tabcolsep}{4pt}
\begin{tabular}{l cc @{\hspace{1em}} ccc c @{\hspace{1em}} ccc c}
\toprule
 & & & \multicolumn{4}{c}{Reference: transportation} & \multicolumn{4}{c}{Reference: metals}\\
Product group & Share & Treated
 & $d_{\text{mach}}^{\ast}$ & $d_{\text{anim}}$ & $d_{\text{met}}$ & Pool
 & $d_{\text{trans}}^{\ast}$ & $d_{\text{mach}}$ & $d_{\text{anim}}$ & Pool\\
\midrule
Transportation          & 0.685 & 0.051 & 0 & 0 & 0 & I  & 1 & 0 & 0 & II\\
Machinery and electrical& 0.072 & 0.667 & 1 & 0 & 0 & II & 0 & 1 & 0 & I\\
Animal products         & 0.051 & 0.018 & 0 & 1 & 0 & I  & 0 & 0 & 1 & I\\
Metals                  & 0.042 & 0.422 & 0 & 0 & 1 & I  & 0 & 0 & 0 & I\\
\bottomrule
\end{tabular}
\caption{The four largest of the 14 product groups in Chang's sample of
$N=1084$ shipments: share of shipments and share treated, and two codings by
dummies with different reference groups. The starred column is the one the Lasso
is assumed to select. Two groups are pooled if their rows agree in every
selected column; post-Lasso then gives them the same fitted propensity score and
outcome trend. With transportation as the reference, it is pooled with every
group whose dummy is not selected. With metals as the reference, transportation
has its own dummy and, once that dummy is selected, its own level. Source: own
calculations based on the replication data of \textcite{chang2020double}; the
names of the groups are inferred from the HS codes in each group.}
\label{tab:pooling}
\end{table}
\begin{ins}

To formalise this, let $Z\in\{1,\dots,C\}$ be a categorical covariate with
shares $\pi_c=\Pr(Z=c)>0$, and suppose, to isolate its role, that $g_0$ and
$\ell_{10}$ depend on $X$ only through $Z$. Any coding of $Z$ by 0/1 columns
together with an intercept spans all functions of $Z$ if it has full rank. A
post-Lasso fit on the selected columns is constant on the groups of categories
that the selected columns do not distinguish: two categories are pooled if their
rows agree in every selected column. Let $\mathcal G_g$ and $\mathcal G_\ell$ be
the resulting partitions of the categories for the two fits, and suppose, as
holds for reference and threshold coding, that the selected columns and the
intercept span all functions that are constant on these pools. In the
population, the refitted propensity score in pool $G$ is then the share of
treated in the pool, $\hat g_G=\E[g_0\mid Z\in G]$.

\begin{proposition}[Pooling bias] % \label{prop:pool}
Under these conditions,
\[
  B_1(\hat g,\hat\ell)=\sum_{G\in\mathcal G_g}\frac{\Pr(Z\in G)}{p_0(1-\hat g_G)}\,
  \operatorname{Cov}\big(g_0,\,\ell_{10}-\hat\ell\;\big|\;Z\in G\big).
\]
If both fits pool the same categories, $\hat\ell$ is constant on each pool and
the covariance is $\operatorname{Cov}(g_0,\ell_{10}\mid Z\in G)$.
\end{proposition}
\begin{proof}
On a pool $G$, $\hat g=\hat g_G$ is constant and
$\E[\hat g_G-g_0\mid Z\in G]=0$. Hence
$\E[(\hat\ell-\ell_{10})(\hat g_G-g_0)\mid Z\in G]
=\operatorname{Cov}(\ell_{10}-\hat\ell,\,g_0\mid Z\in G)$, and $1-\hat g$ is
constant on $G$. Summing over the pools with weights $\Pr(Z\in G)$ in
Lemma~\ref{lem:bias} gives the result. If $\hat\ell$ is constant on $G$, it
drops out of the covariance.
\end{proof}
Proposition~\ref{prop:pool} is the Remark of Section~\ref{sec:nonsparse} applied
pool by pool. When both fits pool the same categories, each pool contributes the
bias of an unadjusted DiD within the pool,
$\operatorname{Cov}(g_0,\ell_{10}\mid Z\in G)/\{\hat g_G(1-\hat g_G)\}$, weighted
by the pool's share of the treated,
$\Pr(Z\in G\mid D=1)=\Pr(Z\in G)\,\hat g_G/p_0$. Pooling is harmless where
selection or the remaining trend is constant within the pool, and harmful where
both vary together. With a single pool, the proposition reduces to the Remark.

\begin{corollary}[Reference category] % \label{cor:reference}
Use reference coding with reference category $b$, and suppose that in both fits
the pool containing $b$ is $\{b\}\cup U$, where the categories in $U$ share a
common level: $g_c=g_U$ and $\ell_{10,c}=\ell_U$ for $c\in U$. Then, with
$w=\pi_b/(\pi_b+\pi_U)$, the pool $\{b\}\cup U$ contributes
\[
  \frac{\pi_b(1-w)}{p_0(1-\hat g_{\{b\}\cup U})}\,(g_b-g_U)(\ell_{10,b}-\ell_U)
\]
to $B_1$. The contribution is zero if the reference category is typical in
selection ($g_b=g_U$) or in trends ($\ell_{10,b}=\ell_U$). It is non-zero only if
the reference category is atypical in both, its sign is the sign of
$(g_b-g_U)(\ell_{10,b}-\ell_U)$, and it grows with the share $\pi_b$ of the
reference category.
\end{corollary}
\begin{proof}
Within the pool, $g_0$ and $\ell_{10}$ take two values, with probability $w$ on
$b$, so
$\operatorname{Cov}(g_0,\ell_{10}\mid Z\in\{b\}\cup U)=w(1-w)(g_b-g_U)(\ell_{10,b}-\ell_U)$.
Multiply by $\Pr(Z\in\{b\}\cup U)=\pi_b/w$.
\end{proof}
The pool $\{b\}\cup U$ arises when neither fit selects the dummies of $U$, and
this is the regime the Lasso is pushed into. With an atypical reference
category, every typical category needs the same non-zero coefficient. It is
cheaper for the Lasso to fit the common level through the unpenalised intercept
and leave these coefficients at zero, which pools the reference category with
the typical ones. Two remarks follow.
\begin{itemize}
\item \emph{Size.} In the regime of \textcite{kolesar2026fragility}, with many
  categories of equal size, $\pi_b$ is of order $1/p$, so a wrong reference
  category costs only $O(1/p)$ in approximation error and satisfies their
  approximate sparsity condition \eqref{eq:sparsity1}. In DML-DiD, the bias is of
  order $\pi_b$, and it is negligible relative to the standard error only if
  $\pi_b\sqrt N\to0$. A wrong reference category therefore matters as soon as its
  share is not small relative to $N^{-1/2}$. This is the typical case in applied
  work, where a common default is to use the largest category as the reference.
\item \emph{Ordered categories.} Suppose $Z$ is ordered, for example in bins, and
  both selection and trends jump at the same bin. With dummy coding, the Lasso
  must select a separate dummy for every bin on one side of the jump, and each
  dummy rests only on the observations in its own bin. With threshold coding,
  $\mathbf 1\{Z\ge c\}$, a single column describes the jump, and it rests on all
  observations. When the bins are too small for their dummies to be selected,
  dummy coding pools bins on both sides of the jump, and by
  Proposition~\ref{prop:pool} the bias approaches that of the unadjusted DiD.
  Threshold coding avoids this.
\end{itemize}
\end{ins}

% =============================================================================
\section{Inserts: Section 5, simulations}\label{sec:sims}

All numbers in this section are computed from the CSV files on the branch
(\texttt{d1\_*.csv}, \texttt{d23\_*.csv}, \texttt{d4\_*.csv} and
\texttt{d4\_*\_forced.csv}). The tables are in \texttt{tables/}; the scripts
\texttt{tables/build\_d23.py} and \texttt{tables/build\_d4.py} rebuild them. They
need \verb|\usepackage{booktabs}|.

\subsection{Practicalities and estimators (pp.~19--20)}

\paragraph{What to fix.}
\begin{itemize}
\item $n$ is the fold size $N/K$ everywhere else, but here it is the number of
  observations in the Lasso. Use $M$, as in Section~3.3.
\item ``The error $\hat e=y-\bar y$ is calculated in a pilot step and the pilot
  residuals in one update'' is hard to parse.
\item ``I follow the code of Chang (2020) and clip'': Chang's code does not
  clip. It drops the observations with $\hat g$ outside $(0.03,0.97)$. Your runs
  clip, so only the wording changes.
\item The list of estimators does not match the tables. Table~1 has two
  versions of DML-DiD and Abadie's estimator, which the list lacks; OR-DiD
  appears only in D4; and $K=5$ is not true for the first panel of Table~1.
\item ``Each design cell is replicated $R=500$ times'': D4 has $R=200$.
  ``replication $r$ uses the $r$-th child of that seed'': no seed has been
  introduced.
\item ``the product term of Lemma ??'' (twice, also in the caption of Table~1):
  Lemma~4.
\end{itemize}
\where{Replace ``Practicalities'' and ``Estimators'' (p.~19 to the top of p.~20).}
\begin{ins}
\paragraph{First steps.} Covariates are standardised on each training sample,
the intercept is never penalised, and all penalties are plug-in with $c=1.1$ and
$\alpha=0.05$. Let $M$ be the number of observations in the regression at hand:
the controls of the auxiliary sample for the outcome nuisance, and the whole
auxiliary sample for the propensity score. The outcome nuisance solves
\[
  \min_\beta\ \frac1M\sum_i(y_i-q_i'\beta)^2+\lambda\sum_j\hat\psi_j|\beta_j|,
  \qquad
  \lambda=\frac{2c}{\sqrt M}\,\Phi^{-1}\Big(1-\frac{\alpha}{2p}\Big),\quad
  \hat\psi_j=\Big(\frac1M\sum_iq_{ij}^2\hat e_i^2\Big)^{1/2}.
\]
The residuals $\hat e_i$ are first $y_i-\bar y$ and then, in one refinement
step, the residuals of the resulting Lasso fit. This is the rule of
\textcite{belloni2012sparse} that Chang suggests, with $\alpha=0.05$ in place of
a level that tends to zero. The propensity score solves
\[
  \min_\beta\ \frac1M\sum_i\big[\log(1+e^{q_i'\beta})-D_iq_i'\beta\big]+\lambda_g\lVert\beta\rVert_1,
  \qquad \lambda_g=\frac{c}{2\sqrt M}\,\Phi^{-1}\Big(1-\frac{\alpha}{2p}\Big),
\]
a plug-in penalty of the form proposed by \textcite{belloni2016post}. I clip the
fitted propensity scores to $[0.03,0.97]$. This is the window Chang's
replication code uses; his code drops the observations outside it instead.

\paragraph{Estimators.}
\begin{enumerate}
\item \emph{DML-DiD, post-Lasso}: the cross-fitted estimator of Definition~2 with
  $K=5$ folds, where each nuisance is refitted without penalty on its selected
  covariates. This is the main specification.
\item \emph{DML-DiD, Lasso (no refit)}: the same, with the penalised fits plugged
  into the score directly, as in Chang's specification.
\item \emph{Oracle}: the score evaluated at the true $(g_0,\ell_{10})$. Its error
  is sampling error alone.
\item \emph{Unadjusted DiD}: the difference in mean outcome changes between
  treated and controls.
\item \emph{Abadie, plug-in Lasso} (D1 only): the estimator (3) with $\hat g$ from
  a logit Lasso fitted on the full sample, without refit and without
  cross-fitting.
\item \emph{OR-DiD} (D4 only): OLS of $Y(1)-Y(0)$ on the full dictionary among the
  controls, imputed for the treated. It requires fewer regressors than controls.
\end{enumerate}

\paragraph{Replications and reporting.} D1 and D2/D3 use $R=500$ replications per
design cell and D4 uses $R=200$. All random numbers derive from one seed: in
each cell, replication $r$ uses the $r$-th child of a seed sequence built from
that seed and the parameters of the cell, so cells with the same data-generating
process share their data. For each estimator I report the bias, the standard
deviation of $\tilde\theta$ across replications (SD), the coverage of the nominal
95\% interval $\tilde\theta\pm1.96\,\widehat{\mathrm{SE}}$, and the ratio of the
average standard error to the SD (SE/SD). Monte Carlo standard errors (MCSE)
measure the simulation uncertainty of a reported number, which shrinks with $R$;
the SD measures the sampling variability of the estimator itself. For a coverage
rate $c$ the MCSE is $\sqrt{c(1-c)/R}$, about 0.010 at nominal coverage for
$R=500$ and 0.015 for $R=200$. For DML-DiD I also report $B_1$, the product term
of Lemma~\ref{lem:bias}, evaluated with the fold-specific fits on an independent
test draw of 5{,}000 observations and averaged over folds.

Each design tests specific results of Section~4. D1 tests Lemma~\ref{lem:bias}
and the Remark of Section~\ref{sec:nonsparse} in a sparse design. D2 and D3 test
Propositions~\ref{prop:ratefail}--\ref{prop:stein} when the coefficients are
dense. D4 tests Lemma~\ref{lem:invariance}, Proposition~\ref{prop:pool} and
Corollary~\ref{cor:reference} when the truth is sparse in one coding but not in
another.
\end{ins}

\subsection{D1 (pp.~20--22)}

\paragraph{What to fix.}
\begin{itemize}
\item \textbf{Table~1 is from the run before your rerun.} The file
  \texttt{d1\_table.tex} in the repository is built from the current CSV files
  and differs in the third decimal in ten entries (for example Abadie
  $1.428\to1.429$, $1.089\to1.090$, $0.944\to0.945$; robust SD of post-Lasso at
  $N=200$ $0.530\to0.524$; coverage of Lasso in the first panel $0.946\to0.948$).
  Nothing changes in substance, but paste the current file so that the thesis
  matches the files you hand in.
\item ``before any notion of sparsity is introduced'': the design is exactly
  sparse. Write ``in a design where sparsity holds''.
\item ``when parallel trends do not hold unconditionally for the untreated
  units'' should be ``when parallel trends hold only conditionally on the
  covariates''.
\item ``$\ell_{10}=1+X'\gamma$ with $\gamma=\delta$'' should be
  $\ell_{10}=1+X'\delta$ with $\delta=\gamma$.
\item ``the same five covariates with the same weights will drive the trend for
  the treated and the untreated'' is not what the design does. The same five
  covariates drive \emph{selection into treatment} and \emph{the untreated
  trend}.
\item ``I a signal-to-noise-ratio such that it explains half the variation in
  the untreated outcome'': a verb is missing, and it is the \emph{change} in the
  outcome.
\item ``as per~.'': the Remark.
\item Results: ``it shows that the plug-in Abadie estimator is skewed'' should be
  ``biased''. ``between $1.563-1.569$'': the table has 1.559--1.566. ``removes
  most of the bias: at $N=200$'': at $N=200$ a third of it remains. ``the
  refitting removes some of the confounding'' says too little: without the refit
  coverage is zero in every trend cell.
\item Digest: ``requieres $N$ to be around 1000 for close to full coverage''
  \textbf{(wrong as printed)}: coverage is 0.884 at $N=1000$ and reaches the
  nominal level only at $N=2000$. ``refitting the penalty'': the penalty is not
  refitted; the coefficients are refitted after selection.
\item The digest says ``D1 shows that the bias of DML-DiD is the product term of
  lemma 4'', but the results text never compares the bias with $B_1$.
\end{itemize}

\where{Replace ``Design setup'' from ``I deviate from his setup \dots'' to ``\dots
while maintaining letting $p=N/2$.'' (p.~20), and add the first paragraph after
the display for $Y(1)-Y(0)$.}
\begin{ins}
The level coefficients $\beta$ cancel in the difference. Whether the outcome
\emph{levels} are sparse is therefore irrelevant for DML-DiD; in the code,
$\beta_j=\gamma_j+0.5$, so every covariate shifts the level. Only the untreated
trend $\ell_{10}$ and the propensity score matter. Following
\textcite{chang2020double}, the first version of the design has
$\ell_{10}\equiv1$, $N=200$, $p=100$, $K=2$, and errors with standard deviation
0.1.

I depart from his design in three ways. First, I use $K=5$ folds. Second, I add
a version with a trend, $\ell_{10}(X)=1+X'\delta$ with $\delta=\gamma$: the five
covariates that drive selection into treatment also drive the untreated trend,
so parallel trends hold conditionally on $X$ but not unconditionally. In this
version the standard deviation of $\varepsilon_2$ is $\|\delta\|=\sqrt{2.2}$, so
that in the population the covariates explain half of the variance of the
untreated change in outcomes. Third, I let $N\in\{200,1000,2000\}$ with
$p=N/2$. With this $\gamma$, the propensity index has standard deviation 1.48, so
overlap is thin in the tails: some controls receive weights $\hat g/(1-\hat g)$
above 30, which inflates the variance of the estimators that use the propensity
score.
\end{ins}
\note{``above 30'' is the clipping bound, $0.97/0.03=32$. Keep the sentence only
if the clipping binds in your D1 runs (check \texttt{np.mean(g == 0.97)} once).}

\where{Replace ``Priors'' (pp.~20--21).}
\begin{ins}
\paragraph{Priors.} In Chang's design the untreated trend is constant, so
$\operatorname{Cov}(g_0,\ell_{10})=0$ and, by the Remark of
Section~\ref{sec:nonsparse}, the unadjusted DiD is unbiased: there is no
confounding to remove. DML-DiD should be unbiased as well. Abadie's estimator is
the score $\psi_1$ with $\ell\equiv0$, so by Lemma~\ref{lem:bias} its bias for a
given propensity fit is
$B_1(\hat g,0)=\E\big[\ell_{10}\,(g_0-\hat g)/\{p_0(1-\hat g)\}\big]$. In Chang's
design, this is the common trend, $\ell_{10}\equiv1$, times the average error in
the weights of the controls, which is not zero when $\hat g$ is shrunk. I
therefore expect Abadie's estimator to be biased, by an amount proportional to
the common trend, although nothing confounds the comparison.

With the trend, the unadjusted DiD is biased by
\[
  \frac{\operatorname{Cov}(g_0,\ell_{10})}{p_0(1-p_0)}
  =\frac{\E[\Lambda'(X'\gamma)]\,\gamma'\delta}{p_0(1-p_0)}=1.563,
\]
since $p_0=0.5$, $\gamma'\delta=2.2$ and $\E[\Lambda'(X'\gamma)]=0.178$. DML-DiD
is consistent, but at small $N$ the smallest coefficients lie below the plug-in
penalty and the Lasso sets them to zero. Lemma~\ref{lem:bias} then predicts a
finite-sample bias equal to $B_1$, which falls as $N$ grows, and
Proposition~\ref{prop:stein} predicts that without the refit the shrinkage of the
selected coefficients adds to it.
\end{ins}

\where{Replace the two paragraphs of results after Table~1 (pp.~21--22).}
\begin{ins}
Table~\ref{tab:d1} confirms each prediction. In Chang's design, DML-DiD is
unbiased (bias below 0.001 in absolute value, MCSE 0.001) and covers at 0.960
($K=2$) and 0.956 ($K=5$). Abadie's plug-in estimator is biased upwards by 0.071
(MCSE 0.002), about 1.7 of its own standard deviations. This reproduces the
contrast Chang uses to motivate DML-DiD. The unadjusted DiD, however, is also
unbiased, covers at 0.948, and is the most precise of the estimators (SD 0.017
against 0.024 for DML-DiD).\footnote{The oracle is less precise than DML-DiD
here (SD 0.027 against 0.024). It uses the true propensity scores, which are
more extreme than the fitted ones and give some controls large weights. That an
estimated propensity score can give a more precise estimator than the true one
is known from \textcite{hirano2003efficient}.} In this design, adjusting for
covariates buys nothing.

Once the covariates enter the trend, the unadjusted DiD is biased by
1.559--1.566 across $N$, matching the predicted 1.563, and never covers.
DML-DiD with post-Lasso first steps removes this bias only gradually. At
$N=200$ the outcome Lasso selects on average 0.21 of the five trend
coefficients, so $\hat\ell_{1k}$ is close to a constant, and the adjustment comes
almost entirely from the propensity score, where the Lasso selects 1.3
covariates. The bias is 0.565 (MCSE 0.028), a third of the unadjusted bias,
coverage is 0.604, and the distribution of $\tilde\theta$ has heavy tails (SD
0.64 against a robust SD of 0.52; SE/SD 0.84). At $N=1000$ and $2000$, on
average 3.07 and 3.95 of the trend coefficients are selected, the bias falls to
0.096 and 0.033, and coverage rises to 0.884 and 0.946 (MCSE 0.014 and 0.010).
Without the refit, as in Chang's specification, the selected coefficients are
shrunk towards zero, and DML-DiD removes much less of the confounding: the bias
is 1.431, 0.776 and 0.490, and coverage is at most 0.004. In every cell, $B_1$
tracks the bias (with the refit 0.617, 0.094 and 0.036 against 0.565, 0.096 and
0.033; without it 1.429, 0.781 and 0.486 against 1.431, 0.776 and 0.490), as
Lemma~\ref{lem:bias} and Lemma~\ref{lem:decomp} require. The oracle covers at
0.946--0.964 throughout, so the shortfall comes from the first steps, not from
the variance estimator. Abadie's plug-in estimator remains biased by 0.95--1.43.
\end{ins}

\where{Replace ``Digest'' (p.~22).}
\begin{ins}
\paragraph{What D1 shows.} First, the bias of DML-DiD is the product term $B_1$
of Lemma~\ref{lem:bias}, in every cell and for both versions of the first steps.
Second, Chang's simulation design contains no confounding: the unadjusted DiD is
unbiased there and more precise than DML-DiD. His simulation shows that a
plug-in first step harms Abadie's estimator and that the orthogonal score
repairs this. It does not show that DML-DiD removes confounding. Third, once the
covariates do confound the trend, DML-DiD needs a large sample even in an
exactly sparse design with five relevant covariates: coverage is 0.60 at
$N=200$ and 0.88 at $N=1000$, and reaches the nominal level at $N=2000$. Fourth,
the refit matters. Without it, the shrinkage of the selected coefficients leaves
half of the confounding at $N=1000$ and a third at $N=2000$, and coverage is
zero.
\end{ins}

\subsection{D2 and D3 (pp.~22--23)}

\paragraph{What to fix in the text that is there.}
\begin{itemize}
\item D3 is never defined. In your code D2 is the design with aligned signs and
  D3 the one with random signs.
\item ``$X\sim N(0,1)$'' should be $X\sim N(0,I_p)$; ``$g_0=\Lambda(X'\delta)$''
  should be $\Lambda(X'\gamma)$ \textbf{(wrong as printed)}; ``let both the outcome
  and the propensity nuisance follow a trend'' is not what is meant.
\item ``$R=200$'': the CSV files have $R=500$.
\item ``the estimates with refitted penalties'': see D1.
\item Priors: for $(2,2)$, ``coverage is expected to approach 1'' should be
  ``approach the nominal level''. The priors contain no numbers, although the
  theory gives them.
\item ``Results'' and ``Digest'' are empty.
\end{itemize}

\where{Replace the whole of ``Second and third design'' (pp.~22--23).}
\begin{ins}
\subsection*{Second and third design: decaying coefficients}

D2 and D3 address the case in which the nuisance functions are not sparse in
the dictionary the researcher uses: every covariate matters a little, as in the
dense models of \textcite{giannone2021economic}. The coefficients decay
polynomially, and the rate of decay moves the design from dense to
approximately sparse. Designs with polynomially decaying coefficients are
standard in this literature
\parencite{belloni2014inference,belloni2017program,athey2018approximate}; the
combination of different decay rates for the two nuisances follows
\textcite[Appendix~D]{smucler2019unifying}.

\paragraph{Design.}
$X_i\sim N(0,I_p)$ with $p=N/2$, $g_0(X)=\Lambda(X'\gamma)$ and
$\ell_{10}(X)=1+X'\delta$, where $|\gamma_j|\propto j^{-\alpha_g}$,
$|\delta_j|\propto j^{-\alpha_\ell}$ and $\|\gamma\|_2=\|\delta\|_2=1$, so that the
strength of selection and of the trend is the same in every cell and only the
decay differs. The outcome is
$Y_i(1)-Y_i(0)=\ell_{10}(X_i)+D_i\theta_0+\varepsilon_i$ with $\theta_0=1$ and
$\varepsilon_i\sim N(0,1)$. In D2 the signs are aligned, $\gamma_j,\delta_j>0$.
In D3 the signs are independent random draws, fixed once per cell and nested
across $N$, so that the leading coefficients are the same for every $N$. The
decay rates are
$(\alpha_g,\alpha_\ell)\in\{(0.7,0.7),(2,2),(0.8,1.5)\}$,
$N\in\{500,1000,2000,4000\}$, and $K=5$.

\paragraph{Priors.} With coefficients that decay like $j^{-\alpha}$, the Lasso
error is of order $N^{-(1-e)/2}$ with $e=1/(2\alpha)$, up to logarithms
\parencite[Example~10]{smucler2019unifying}. Chang's condition, an error below
$N^{-1/4}$ for each first step, holds if $\alpha>1$ for both. The product
condition of Proposition~\ref{prop:sparse} holds if
$1/\alpha_g+1/\alpha_\ell<2$. The three cells are chosen accordingly:
$(0.7,0.7)$ violates both conditions, $(2,2)$ satisfies both, and $(0.8,1.5)$
satisfies the product condition but not Chang's. Four predictions follow.
\begin{enumerate}
\item By Lemma~\ref{lem:bias}, the bias equals $B_1$ in every cell, and by
  Proposition~\ref{prop:stein}, $B_1$ is given by the errors of the two
  coefficient vectors.
\item The bias relative to the standard error, $\sqrt N\,|\text{bias}|$, behaves
  like $N^{(e_g+e_\ell-1)/2}$. It grows in $(0.7,0.7)$, with exponent $0.21$,
  so coverage goes to zero (Proposition~\ref{prop:ratefail}); it falls in
  $(2,2)$, with exponent $-0.25$, so coverage approaches the nominal level
  (Proposition~\ref{prop:ratedr}); and in $(0.8,1.5)$ the exponent is $-0.02$:
  the product condition holds, but so narrowly that the bias should stay
  roughly constant relative to the standard error at any realistic $N$.
\item By \eqref{eq:both}, the bias is the sum of $\gamma_j\delta_j$ over the
  covariates that both fits omit. With aligned signs every term is positive;
  with random signs the terms partly cancel. The bias should therefore be
  smaller in D3 than in D2 wherever many covariates are omitted.
\item Without the refit, the shrinkage of the selected coefficients adds to the
  bias in every cell.
\end{enumerate}
As a benchmark, the bias of the unadjusted DiD is $0.826\,\gamma'\delta$ in all
cells, by the Remark of Section~\ref{sec:nonsparse} with $p_0=0.5$ and
$\E[\Lambda'(X'\gamma)]=0.207$.

\end{ins}
\begin{table}[ht]
\centering\footnotesize
\setlength{\tabcolsep}{3pt}
\begin{tabular}{ll rr rr rr rr rr r r}
\toprule
 & & & & \multicolumn{2}{c}{$N=500$} & \multicolumn{2}{c}{$N=1000$} & \multicolumn{2}{c}{$N=2000$} & \multicolumn{2}{c}{$N=4000$} & \multicolumn{2}{c}{Slope} \\
\cmidrule(lr){5-6} \cmidrule(lr){7-8} \cmidrule(lr){9-10} \cmidrule(lr){11-12} \cmidrule(lr){13-14}
$(\alpha_g,\alpha_\ell)$ & Signs & $\gamma'\delta$ & Unadj. & Bias & Cov. & Bias & Cov. & Bias & Cov. & Bias & Cov. & Est. & Pred. \\
\midrule
\multicolumn{14}{l}{\emph{A. DML-DiD, post-Lasso}} \\
$(0.7,\,0.7)$ & aligned & 1.00 & 0.83 & 0.511 & 0.04 & 0.445 & 0.00 & 0.377 & 0.00 & 0.320 & 0.00 & 0.27 & 0.21 \\
$(0.7,\,0.7)$ & random & $-$0.52 & $-$0.43 & $-$0.146 & 0.78 & $-$0.091 & 0.80 & $-$0.079 & 0.72 & $-$0.058 & 0.73 & -- & -- \\
$(2,\,2)$ & aligned & 1.00 & 0.83 & 0.065 & 0.93 & 0.045 & 0.90 & 0.018 & 0.95 & 0.016 & 0.93 & $-$0.25 & $-$0.25 \\
$(2,\,2)$ & random & $-$0.85 & $-$0.71 & 0.073 & 0.91 & 0.036 & 0.91 & 0.014 & 0.94 & 0.010 & 0.96 & -- & -- \\
$(0.8,\,1.5)$ & aligned & 0.87 & 0.72 & 0.180 & 0.61 & 0.121 & 0.66 & 0.083 & 0.66 & 0.057 & 0.66 & $-$0.05 & $-$0.02 \\
$(0.8,\,1.5)$ & random & 0.76 & 0.63 & 0.095 & 0.84 & 0.022 & 0.93 & 0.001 & 0.95 & 0.000 & 0.95 & -- & -- \\
\addlinespace
\multicolumn{14}{l}{\emph{B. DML-DiD, Lasso (no refit)}} \\
$(0.7,\,0.7)$ & aligned & 1.00 & 0.83 & 0.729 & 0.00 & 0.648 & 0.00 & 0.566 & 0.00 & 0.488 & 0.00 & 0.31 & 0.21 \\
$(0.7,\,0.7)$ & random & $-$0.52 & $-$0.43 & $-$0.354 & 0.17 & $-$0.243 & 0.18 & $-$0.173 & 0.15 & $-$0.128 & 0.10 & -- & -- \\
$(2,\,2)$ & aligned & 1.00 & 0.83 & 0.279 & 0.23 & 0.181 & 0.29 & 0.121 & 0.31 & 0.079 & 0.41 & $-$0.11 & $-$0.25 \\
$(2,\,2)$ & random & $-$0.85 & $-$0.71 & $-$0.133 & 0.72 & $-$0.044 & 0.87 & $-$0.009 & 0.91 & 0.003 & 0.94 & -- & -- \\
$(0.8,\,1.5)$ & aligned & 0.87 & 0.72 & 0.402 & 0.02 & 0.301 & 0.00 & 0.220 & 0.00 & 0.161 & 0.00 & 0.06 & $-$0.02 \\
$(0.8,\,1.5)$ & random & 0.76 & 0.63 & 0.316 & 0.15 & 0.204 & 0.17 & 0.123 & 0.28 & 0.067 & 0.49 & -- & -- \\
\bottomrule
\end{tabular}

\caption{D2 (aligned signs) and D3 (random signs): bias and coverage of the
nominal 95\% interval of DML-DiD, $R=500$, $\theta_0=1$. $\gamma'\delta$ and the
bias of the unadjusted DiD are given for $N=4000$. ``Slope'' is the slope of
$\log(\sqrt N|\text{bias}|)$ on $\log N$ and its predicted value
$(e_g+e_\ell-1)/2$. The largest MCSE of a bias in the table is 0.006.}
\label{tab:d23}
\end{table}
\begin{ins}

\end{ins}
\begin{table}[ht]
\centering\footnotesize
\setlength{\tabcolsep}{4pt}
\begin{tabular}{ll r rrrr rrrr}
\toprule
 & & & \multicolumn{4}{c}{Post-Lasso} & \multicolumn{4}{c}{Lasso (no refit)} \\
\cmidrule(lr){4-7} \cmidrule(lr){8-11}
$(\alpha_g,\alpha_\ell)$ & Signs & $N$ & Bias & $B_1$ & Formula & Both omit & Bias & $B_1$ & Inner & Both omit \\
\midrule
$(0.7,\,0.7)$ & aligned & 500 & 0.511 & 0.528 & 0.520 & 0.655 & 0.729 & 0.734 & 0.900 & 0.630 \\
$(0.7,\,0.7)$ & aligned & 4000 & 0.320 & 0.320 & 0.317 & 0.357 & 0.488 & 0.486 & 0.543 & 0.336 \\
$(0.7,\,0.7)$ & random & 500 & $-$0.146 & $-$0.138 & $-$0.139 & $-$0.176 & $-$0.354 & $-$0.354 & $-$0.418 & $-$0.171 \\
$(0.7,\,0.7)$ & random & 4000 & $-$0.058 & $-$0.057 & $-$0.057 & $-$0.063 & $-$0.128 & $-$0.124 & $-$0.147 & $-$0.060 \\
\addlinespace
$(2,\,2)$ & aligned & 500 & 0.065 & 0.074 & 0.070 & 0.076 & 0.279 & 0.280 & 0.307 & 0.069 \\
$(2,\,2)$ & aligned & 4000 & 0.016 & 0.014 & 0.014 & 0.015 & 0.079 & 0.077 & 0.076 & 0.014 \\
$(2,\,2)$ & random & 500 & 0.073 & 0.062 & 0.060 & 0.067 & $-$0.133 & $-$0.137 & $-$0.152 & 0.060 \\
$(2,\,2)$ & random & 4000 & 0.010 & 0.008 & 0.008 & 0.009 & 0.003 & 0.004 & 0.003 & 0.009 \\
\addlinespace
$(0.8,\,1.5)$ & aligned & 500 & 0.180 & 0.197 & 0.190 & 0.240 & 0.402 & 0.408 & 0.498 & 0.230 \\
$(0.8,\,1.5)$ & aligned & 4000 & 0.057 & 0.055 & 0.055 & 0.064 & 0.161 & 0.158 & 0.177 & 0.060 \\
$(0.8,\,1.5)$ & random & 500 & 0.095 & 0.088 & 0.091 & 0.123 & 0.316 & 0.311 & 0.377 & 0.116 \\
$(0.8,\,1.5)$ & random & 4000 & 0.000 & 0.001 & 0.002 & 0.001 & 0.067 & 0.069 & 0.078 & 0.001 \\
\bottomrule
\end{tabular}

\caption{D2 and D3: bias of DML-DiD, $B_1$ on the test draw, and the
approximations of Section~\ref{sec:nonsparse}. ``Formula'' is
Proposition~\ref{prop:stein} evaluated at the fitted coefficients, ``Inner'' is
\eqref{eq:inner} and ``Both omit'' is \eqref{eq:both}.}
\label{tab:d23formula}
\end{table}
\begin{ins}

\paragraph{Results.} Tables~\ref{tab:d23} and~\ref{tab:d23formula} show the
following.
\begin{enumerate}
\item \emph{The bias is $B_1$.} In every cell the bias of DML-DiD and $B_1$ agree
  to within 0.02 at $N=500$ and 0.004 at $N=4000$, and
  Proposition~\ref{prop:stein} reproduces $B_1$ to within 0.01 for post-Lasso.
  The standard errors themselves are accurate: SE/SD is between 0.92 and 1.05
  in all cells. The failures of coverage below are failures of centring, not of
  the variance estimator.
\item \emph{Dense in both nuisances.} In $(0.7,0.7)$ with aligned signs,
  post-Lasso leaves 62\% of the bias of the unadjusted DiD at $N=500$ and 39\% at
  $N=4000$, where the Lasso selects on average 5 of the 2{,}000 covariates for
  the propensity score and 6 for the trend. $\sqrt N\,|\text{bias}|$ grows from
  11 to 20 (slope 0.27, against a predicted 0.21), and coverage is 0.04 at
  $N=500$ and zero thereafter, as Proposition~\ref{prop:ratefail} predicts.
\item \emph{Sparse in both.} In $(2,2)$, the bias of post-Lasso falls from 0.065
  to 0.016, the slope of $\sqrt N\,|\text{bias}|$ is $-0.25$, as predicted, and
  coverage is between 0.90 and 0.95.
\item \emph{Inside the product region, outside Chang's condition.} In
  $(0.8,1.5)$ with aligned signs, $\sqrt N\,|\text{bias}|$ barely moves, from 4.0
  at $N=500$ to 3.6 at $N=4000$ (slope $-0.05$, against a predicted $-0.02$), and
  coverage stays at 0.61--0.66. Proposition~\ref{prop:ratedr} applies to this
  cell, but its conclusion is not visible at these sample sizes: an exponent of
  $-0.02$ would require $N$ to grow by many orders of magnitude before the bias
  becomes small relative to the standard error.
\item \emph{Alignment.} With random signs, the contributions of the omitted
  covariates partly cancel. In $(0.8,1.5)$ the bias of post-Lasso is 0.022 at
  $N=1000$ and vanishes from $N=2000$, with coverage 0.93--0.95, against 0.66
  with aligned signs. In $(0.7,0.7)$ the share of the unadjusted bias that
  remains at $N=4000$ falls from 39\% to 13\%, but coverage is still only
  0.72--0.80. In $(2,2)$, where few covariates are omitted, alignment makes
  little difference. For random signs I do not report slopes: the bias falls
  below its Monte Carlo error or changes sign as more covariates are selected,
  so $\log|\text{bias}|$ does not measure a rate.
\item \emph{The refit.} Without the refit, the bias is 1.4 to 7 times larger in
  the cells with aligned signs, and coverage is at most 0.41, even in $(2,2)$.
  For post-Lasso, the sum over the covariates that both fits omit,
  \eqref{eq:both}, tracks $B_1$ and over-predicts it by 3--24\%. Without the
  refit it misses the shrinkage of the selected coefficients (0.069 against a
  $B_1$ of 0.280 in $(2,2)$ at $N=500$), and only the inner product
  \eqref{eq:inner} tracks $B_1$.
\end{enumerate}

\paragraph{What D2 and D3 show.} When both nuisances are dense in the
researcher's dictionary, DML-DiD with Lasso first steps removes only part of the
confounding and its confidence interval does not cover, although the standard
errors are accurate. The product condition describes the limit correctly, but
close to its boundary the limit is irrelevant in practice: a design that
satisfies the condition has a coverage of two thirds at every sample size
considered. Whether the remaining bias is large depends on whether the omitted
covariates push selection and trends in the same direction, which the
researcher cannot check from the selected model. The refit matters in every
cell.
\end{ins}
\note{$\E[\Lambda'(X'\gamma)]=0.207$ is implied by your notebook output
(unadjusted bias $0.826$ at $\gamma'\delta=1$, $p_0=0.5$); print
\texttt{e\_dlam} once to confirm the third decimal. If your supervisor wants
D2/D3 shorter, move Table~\ref{tab:d23formula} and result~6 to the appendix; the
main text then needs only Table~\ref{tab:d23} with the aligned rows
(\texttt{tables/d23\_table\_main.tex}) and results 1--5.}

\subsection{D4 (pp.~23--24)}

\paragraph{What to fix in the text that is there.}
\begin{itemize}
\item ``investigates the how DML-DiD grapples in the setting of Kolesár et al.''
  needs rewriting; ``categoricals and dummies are included in the structural
  model, where only one of the representations of the covariate matrix is truly
  sparse'' is not accurate: several codings are sparse, and one is not.
\item The setup stops in the middle of a sentence (``Category''), numbers the
  categories 1--3 where your code and my tables use 0--2, and does not give
  $C$, the shares, the codings, D4b, $N$, $R$ or the estimators.
\item ``By corollary ??'' is Corollary~14; ``By Lemma ??'' is Lemma~10. There is
  a stray ``I'' on p.~24.
\item The ``Digest'' is my list of priors from the chat, pasted as running text
  (``P10, \dots\ predicted at 0.045, 0.119 and 0.195 for $=0.05$ \dots\ equals B
  and the bias''). It is not a digest and several symbols were lost.
\end{itemize}

\where{Replace the whole of ``Fourth design'' (pp.~23--24).}
\begin{ins}
\subsection*{Fourth design: the same truth in different codings}

D4 addresses the case in which the nuisance functions are sparse in one coding
of the covariates but not in another. It transfers the categorical examples of
\textcite[Section~3.3]{kolesar2026fragility} to DML-DiD. The truth is fixed and
only its coding changes: each replication draws one data set and one fold
split, and every coding is estimated on the same data with the same folds, so
differences between codings are pure effects of the representation.

\paragraph{Design.}
A categorical covariate $Z$ with $C$ categories enters both nuisances,
$g_0=\Lambda(-0.25+\gamma_Z)$ and $\ell_{10}=1+\delta_Z$, and
$Y_i(1)-Y_i(0)=\ell_{10}+\theta_0D_i+\varepsilon_i$ with $\theta_0=1$ and
$\varepsilon_i\sim N(0,1)$. Every dictionary also contains $N/10$ irrelevant
$N(0,1)$ covariates, $N\in\{1000,2000,4000\}$, $R=200$ and $K=5$.
\begin{itemize}
\item \emph{D4a (reference category).} $C=20$. Category 0 is atypical in both
  selection and trends, $(\gamma_0,\delta_0)=(1,1)$; category 1 only in
  selection, $(1,0)$; category 2 only in trends, $(0,1)$; categories 3--19 are
  typical, $(0,0)$. Category 0 has share $\pi_0\in\{0.05,0.15,0.30\}$ and the
  other categories share the rest equally. The codings are dummies with
  reference category 3, 1, 2 or 0.
\item \emph{D4b (ordered categories).} $C=40$ ordered categories with equal
  shares, like age bins. Selection and trends jump in the middle,
  $\gamma_c=\mathbf 1\{c\ge20\}$ and $\delta_c=0.75\cdot\mathbf 1\{c\ge20\}$. The
  codings are thresholds, $\mathbf 1\{Z\ge c\}$ for $c=1,\dots,39$, and dummies
  with the lowest category as reference.
\end{itemize}
Besides DML-DiD with and without the refit, I report two estimators that do not
penalise the dummies: OR-DiD on the full dictionary, and DML-DiD with the
dummies \emph{forced in}, where the Lasso penalises only the irrelevant
covariates.

\paragraph{Priors.}
\begin{enumerate}
\item \emph{Reference category.} By Corollary~\ref{cor:reference}, only
  reference category 0, which is atypical in both nuisances, creates bias. If
  the Lasso gives categories 1 and 2 their own levels and pools category 0 with
  the typical ones, the corollary predicts a bias of 0.045, 0.119 and 0.195 for
  $\pi_0=0.05$, 0.15 and 0.30. This exceeds the bias of the unadjusted DiD
  (0.039, 0.108 and 0.180), because categories 1 and 2 contribute to the
  unadjusted bias with the opposite sign. References 1, 2 and 3 are atypical in
  at most one nuisance and should be harmless.
\item \emph{Pooling.} By Proposition~\ref{prop:pool}, the bias equals the pooling
  formula evaluated at the selected columns, in every cell.
\item \emph{Ordered categories.} Threshold coding needs one selected column and
  should be unbiased. Dummy coding needs 20 dummies, each resting on 2.5\% of the
  sample. As long as the Lasso cannot select them, it pools across the jump, and
  the bias approaches that of the unadjusted DiD, 0.184.
\item \emph{Unpenalised dummies.} By Lemma~\ref{lem:invariance}, OR-DiD gives the
  same estimate under every coding and is unbiased, at the price of a larger
  variance, since it estimates all $N/10+C-1$ coefficients. Forcing the dummies
  in makes DML-DiD invariant to the reference category as well. I expect it to
  be unbiased, with a variance between that of DML-DiD with a harmless coding
  and that of OR-DiD.
\item \emph{The range grows with $N$.} Where a coding is biased, the bias stays
  while the standard error shrinks, so the studentised range across codings,
  $\mathcal R(\mathcal A)$, should grow with $N$.
\item \emph{The refit.} Without the refit, the shrinkage of the selected
  coefficients adds bias, also for the harmless codings.
\end{enumerate}

\end{ins}
\begin{table}[ht]
\centering\footnotesize
\setlength{\tabcolsep}{2.8pt}
\begin{tabular}{l rrrrr rr rr rr}
\toprule
 & \multicolumn{5}{c}{Post-Lasso} & & & \multicolumn{2}{c}{No refit} & \multicolumn{2}{c}{Selected} \\
\cmidrule(lr){2-6} \cmidrule(lr){9-10} \cmidrule(lr){11-12}
 & Bias & MCSE & SD & SE/SD & Cov. & $B_1$ & Pool & Bias & Cov. & $\hat g$ & $\hat\ell$ \\
\midrule
\multicolumn{12}{l}{\emph{$\pi_0=0.05$: unadjusted bias 0.039, Corollary for reference 0: 0.045}} \\
\quad Reference 3 (typical) & $-$0.002 & 0.002 & 0.031 & 1.07 & 0.97 & $-$0.001 & 0.000 & 0.025 & 0.91 & 1.9 & 1.9 \\
\quad Reference 1 (selection only) & $-$0.002 & 0.002 & 0.030 & 1.08 & 0.96 & $-$0.001 & 0.000 & 0.024 & 0.92 & 0.9 & 1.9 \\
\quad Reference 2 (trend only) & $-$0.001 & 0.002 & 0.031 & 1.08 & 0.97 & $-$0.001 & 0.000 & 0.024 & 0.93 & 1.9 & 0.9 \\
\quad Reference 0 (both) & 0.044 & 0.002 & 0.031 & 1.06 & 0.77 & 0.046 & 0.045 & 0.041 & 0.76 & 0.9 & 1.0 \\
\quad Dummies forced in & $-$0.002 & 0.002 & 0.031 & 1.07 & 0.97 & 0.000 & -- & $-$0.002 & 0.97 & all & all \\
\quad OR-DiD & $-$0.001 & 0.002 & 0.035 & -- & -- & -- & -- & -- & -- & -- & -- \\
\quad Unadjusted DiD & 0.039 & 0.002 & 0.031 & 1.07 & 0.81 & -- & -- & -- & -- & -- & -- \\
\addlinespace
\multicolumn{12}{l}{\emph{$\pi_0=0.15$: unadjusted bias 0.108, Corollary for reference 0: 0.119}} \\
\quad Reference 3 (typical) & $-$0.001 & 0.002 & 0.032 & 1.06 & 0.98 & 0.000 & 0.000 & 0.023 & 0.90 & 1.8 & 2.0 \\
\quad Reference 1 (selection only) & $-$0.001 & 0.002 & 0.031 & 1.07 & 0.97 & 0.000 & 0.000 & 0.023 & 0.90 & 1.0 & 2.0 \\
\quad Reference 2 (trend only) & $-$0.001 & 0.002 & 0.032 & 1.05 & 0.97 & 0.000 & 0.000 & 0.022 & 0.91 & 1.8 & 1.0 \\
\quad Reference 0 (both) & 0.117 & 0.002 & 0.033 & 1.04 & 0.06 & 0.115 & 0.118 & 0.110 & 0.09 & 0.7 & 1.0 \\
\quad Dummies forced in & $-$0.002 & 0.002 & 0.032 & 1.05 & 0.99 & 0.000 & -- & $-$0.002 & 0.99 & all & all \\
\quad OR-DiD & 0.000 & 0.003 & 0.037 & -- & -- & -- & -- & -- & -- & -- & -- \\
\quad Unadjusted DiD & 0.107 & 0.002 & 0.032 & 1.05 & 0.10 & -- & -- & -- & -- & -- & -- \\
\addlinespace
\multicolumn{12}{l}{\emph{$\pi_0=0.30$: unadjusted bias 0.180, Corollary for reference 0: 0.195}} \\
\quad Reference 3 (typical) & 0.000 & 0.002 & 0.031 & 1.08 & 0.97 & 0.000 & 0.000 & 0.022 & 0.93 & 1.7 & 2.0 \\
\quad Reference 1 (selection only) & 0.000 & 0.002 & 0.030 & 1.11 & 0.98 & 0.000 & 0.000 & 0.022 & 0.93 & 1.0 & 2.0 \\
\quad Reference 2 (trend only) & 0.000 & 0.002 & 0.032 & 1.09 & 0.97 & 0.000 & $-$0.001 & 0.021 & 0.92 & 1.7 & 1.0 \\
\quad Reference 0 (both) & 0.189 & 0.002 & 0.033 & 1.04 & 0.00 & 0.190 & 0.190 & 0.181 & 0.00 & 0.4 & 0.9 \\
\quad Dummies forced in & 0.000 & 0.002 & 0.032 & 1.09 & 0.97 & 0.000 & -- & 0.000 & 0.96 & all & all \\
\quad OR-DiD & 0.001 & 0.003 & 0.038 & -- & -- & -- & -- & -- & -- & -- & -- \\
\quad Unadjusted DiD & 0.179 & 0.002 & 0.033 & 1.04 & 0.00 & -- & -- & -- & -- & -- & -- \\
\bottomrule
\end{tabular}

\caption{D4a at $N=4000$, $R=200$, $\theta_0=1$. The label of each coding says in
which nuisance its reference category is atypical. ``Pool'' is the formula of
Proposition~\ref{prop:pool} evaluated at the selected columns; ``No refit'' is
the same Lasso fit plugged in without refitting; ``Selected'' is the average
number of dummies the Lasso selects for the propensity score and the trend.
OR-DiD is identical under all four reference categories.}
\label{tab:d4a}
\end{table}
\begin{ins}

\end{ins}
\begin{table}[ht]
\centering\small
\begin{tabular}{ll rr rr rr r}
\toprule
 & & \multicolumn{2}{c}{Reference 3} & \multicolumn{2}{c}{Reference 0} & \multicolumn{2}{c}{Forced in} & \\
\cmidrule(lr){3-4} \cmidrule(lr){5-6} \cmidrule(lr){7-8}
$\pi_0$ & $N$ & Bias & Cov. & Bias & Cov. & Bias & Cov. & Range \\
\midrule
0.05 & 1000 & 0.024 & 0.94 & 0.038 & 0.91 & $-$0.002 & 0.95 & 0.12 \\
 & 2000 & 0.007 & 0.94 & 0.043 & 0.83 & 0.000 & 0.94 & 0.86 \\
 & 4000 & $-$0.002 & 0.97 & 0.044 & 0.77 & $-$0.002 & 0.97 & 1.47 \\
\addlinespace
0.15 & 1000 & $-$0.015 & 0.93 & 0.106 & 0.66 & $-$0.005 & 0.94 & 1.75 \\
 & 2000 & $-$0.001 & 0.94 & 0.115 & 0.38 & $-$0.002 & 0.94 & 2.51 \\
 & 4000 & $-$0.001 & 0.98 & 0.117 & 0.06 & $-$0.002 & 0.99 & 3.63 \\
\addlinespace
0.30 & 1000 & $-$0.007 & 0.94 & 0.175 & 0.27 & $-$0.003 & 0.94 & 2.68 \\
 & 2000 & $-$0.003 & 0.96 & 0.182 & 0.04 & $-$0.003 & 0.96 & 3.87 \\
 & 4000 & 0.000 & 0.97 & 0.189 & 0.00 & 0.000 & 0.97 & 5.64 \\
\bottomrule
\end{tabular}

\caption{D4a across $N$: bias and coverage of DML-DiD (post-Lasso) with a typical
reference category, with the reference category that is atypical in both
nuisances, and with the dummies forced in. ``Range'' is the median over
replications of the studentised range across the four reference categories,
in standard errors of the estimate with reference 3.}
\label{tab:d4aN}
\end{table}
\begin{ins}

\end{ins}
\begin{table}[ht]
\centering\footnotesize
\setlength{\tabcolsep}{3.5pt}
\begin{tabular}{l rrrrr rr rr rr}
\toprule
 & \multicolumn{5}{c}{Post-Lasso} & & & \multicolumn{2}{c}{No refit} & \multicolumn{2}{c}{Selected} \\
\cmidrule(lr){2-6} \cmidrule(lr){9-10} \cmidrule(lr){11-12}
 & Bias & MCSE & SD & SE/SD & Cov. & $B_1$ & Pool & Bias & Cov. & $\hat g$ & $\hat\ell$ \\
\midrule
\multicolumn{12}{l}{\emph{$N=1000$: median studentised range 2.58}} \\
\quad Thresholds & $-$0.005 & 0.005 & 0.070 & 1.00 & 0.94 & 0.001 & 0.000 & 0.056 & 0.85 & 1.7 & 1.7 \\
\quad Dummies, reference 0 & 0.175 & 0.005 & 0.069 & 0.99 & 0.28 & 0.183 & 0.183 & 0.176 & 0.28 & 0.0 & 0.0 \\
\quad Dummies forced in & $-$0.014 & 0.009 & 0.122 & 0.88 & 0.96 & $-$0.009 & -- & $-$0.008 & 0.93 & all & all \\
\quad OR-DiD & $-$0.002 & 0.006 & 0.079 & -- & -- & -- & -- & -- & -- & -- & -- \\
\quad Unadjusted DiD & 0.175 & 0.005 & 0.069 & 0.98 & 0.26 & -- & -- & -- & -- & -- & -- \\
\addlinespace
\multicolumn{12}{l}{\emph{$N=2000$: median studentised range 3.70}} \\
\quad Thresholds & 0.007 & 0.003 & 0.047 & 1.04 & 0.95 & 0.001 & 0.000 & 0.042 & 0.87 & 2.0 & 1.9 \\
\quad Dummies, reference 0 & 0.190 & 0.003 & 0.045 & 1.10 & 0.03 & 0.180 & 0.183 & 0.192 & 0.01 & 0.1 & 0.0 \\
\quad Dummies forced in & 0.006 & 0.004 & 0.051 & 1.11 & 0.98 & $-$0.001 & -- & 0.006 & 0.97 & all & all \\
\quad OR-DiD & 0.006 & 0.004 & 0.053 & -- & -- & -- & -- & -- & -- & -- & -- \\
\quad Unadjusted DiD & 0.192 & 0.003 & 0.043 & 1.11 & 0.01 & -- & -- & -- & -- & -- & -- \\
\addlinespace
\multicolumn{12}{l}{\emph{$N=4000$: median studentised range 5.06}} \\
\quad Thresholds & 0.000 & 0.002 & 0.033 & 1.03 & 0.95 & 0.000 & 0.000 & 0.019 & 0.93 & 1.9 & 1.9 \\
\quad Dummies, reference 0 & 0.173 & 0.002 & 0.032 & 1.07 & 0.00 & 0.176 & 0.181 & 0.180 & 0.00 & 0.6 & 0.1 \\
\quad Dummies forced in & 0.000 & 0.002 & 0.035 & 1.04 & 0.95 & 0.000 & -- & 0.000 & 0.94 & all & all \\
\quad OR-DiD & $-$0.001 & 0.003 & 0.038 & -- & -- & -- & -- & -- & -- & -- & -- \\
\quad Unadjusted DiD & 0.181 & 0.002 & 0.032 & 1.07 & 0.00 & -- & -- & -- & -- & -- & -- \\
\bottomrule
\end{tabular}

\caption{D4b: ordered categories with a jump in the middle, $R=200$,
$\theta_0=1$. The bias of the unadjusted DiD is 0.184 in the population. Columns
as in Table~\ref{tab:d4a}.}
\label{tab:d4b}
\end{table}
\begin{ins}

\paragraph{Results.} Tables~\ref{tab:d4a}--\ref{tab:d4b} confirm the priors.
\begin{enumerate}
\item \emph{Only the reference category that is atypical in both nuisances
  creates bias.} At $N=4000$ its bias is 0.044, 0.117 and 0.189 for
  $\pi_0=0.05$, 0.15 and 0.30 (MCSE 0.002), against 0.045, 0.119 and 0.195
  predicted by Corollary~\ref{cor:reference}, and coverage is 0.77, 0.06 and
  0.00. With references 1, 2 and 3 the bias is at most 0.002 in absolute value
  and coverage is 0.96--0.98. A reference category that is atypical in only one
  nuisance is harmless: this is the categorical form of the product structure
  of Lemma~\ref{lem:bias}.
\item \emph{The pooling formula gives the bias.} $B_1$ agrees with the bias to
  within 0.010 in every cell, and to within 0.003 in D4a at $N=4000$. The formula
  of Proposition~\ref{prop:pool}, evaluated at the selected columns, is a
  population quantity that ignores the sampling error of the refit; it agrees
  with the bias to within 0.015, with the largest gaps at $N=1000$.
\item \emph{The bias does not shrink with $N$.} For $\pi_0=0.30$ the bias of
  reference 0 is 0.175, 0.182 and 0.189 at $N=1000$, 2000 and 4000, while its
  standard error falls, so coverage drops from 0.27 to 0.00
  (Table~\ref{tab:d4aN}). At $N=4000$ the Lasso selects on average 0.4 and 0.9
  of the 19 dummies. With $C$ fixed, the dummies of the typical categories
  would eventually be selected, but not at these sample sizes.
\item \emph{A bad reference category is worse than no adjustment.} With
  reference 0 the bias exceeds that of the unadjusted DiD (0.189 against 0.179
  for $\pi_0=0.30$, and likewise for the other shares), because the Lasso does
  fit categories 1 and 2, whose contributions to the unadjusted bias have the
  opposite sign.
\item \emph{When the Lasso selects almost nothing, the coding does not matter.}
  At $N=1000$ and $\pi_0=0.05$, every coding pools nearly all categories, all
  codings are biased (0.024 with reference 3 and 0.038 with reference 0), and
  the median studentised range is 0.12. A small range is therefore not evidence
  that a specification is unbiased.
\item \emph{Ordered categories.} In D4b, threshold coding is unbiased (bias at
  most 0.007 in absolute value, coverage 0.94--0.95). Dummy coding stays at the
  bias of the unadjusted DiD (0.175, 0.190 and 0.173, against 0.184), with
  coverage 0.28, 0.03 and 0.00.
\item \emph{Unpenalised dummies.} OR-DiD gives the same estimate under every
  coding, to machine precision, and is unbiased, with a standard deviation 13--23\%
  larger than that of DML-DiD with a harmless coding at $N=4000$. With the
  dummies forced in, DML-DiD is unbiased in D4a (bias at most 0.005 in absolute
  value, coverage 0.94--0.99), and its standard deviation is within 6\% of that
  of DML-DiD with a harmless coding. Forcing in is not free when the categories
  are small: in D4b at $N=1000$, where each of the 40 bins holds 25
  observations, the standard deviation is 0.122 against 0.070 for threshold
  coding, and the standard error understates it (SE/SD 0.88). At $N=2000$ and
  4000 the difference is 9\% and 6\%.
\item \emph{The range grows with $N$.} The median studentised range across the
  four reference categories grows from 0.12, 1.75 and 2.68 at $N=1000$ to 1.47,
  3.63 and 5.64 at $N=4000$ for the three shares, and from 2.58 to 5.06 across
  the two codings of D4b. This is the order of magnitude that
  \textcite{kolesar2026fragility} document in their applications, where such
  choices move sparsity-based estimates by two standard errors or more.
\item \emph{The refit.} Without the refit, the harmless codings are biased by
  0.02--0.07, with coverage 0.82--0.93.
\end{enumerate}

\paragraph{What D4 shows.} The coding of a categorical covariate matters for
DML-DiD exactly where Proposition~\ref{prop:pool} says it should: when the Lasso
pools categories that differ in both selection and trends. The damage grows with
the share of the pooled category, it does not fade over the range of $N$
considered, and it can exceed the bias of making no adjustment at all. A
researcher can avoid the choice by not penalising the dummies, which costs
little unless the categories are very small, or, for ordered categories, by
threshold coding. D4 uses a single categorical covariate with a known structure.
Section~\ref{sec:emp} asks how much these choices matter in Chang's application.
\end{ins}
\note{(i) The forced-in files contain only \texttt{forced\_ref3} (and
\texttt{forced\_ref0} for D4b). The invariance of the forced-in estimate to the
reference category is therefore a consequence of Lemma~10, not something the
tables show. If you want to show it, run one cell with a second reference; since
your code implements ``unpenalised'' as a penalty divided by 100, the two
estimates will agree to about four decimals, not exactly, and the text should
say so. (ii) The three tables take most of two pages. If
space is short, keep Tables~\ref{tab:d4aN} and~\ref{tab:d4b} in the main text
and move Table~\ref{tab:d4a} to the appendix.}

% =============================================================================
\section{Inserts: Section 6, Chang's application}\label{sec:emp-ins}

\subsection{State of the section}
Section~6 of v13 (pp.~24--25) is still the plan from v11: an opening paragraph
in which five citations are broken, the summary of Sequeira's paper, a paragraph
of intentions (``I intend to conflate the covariate matrix \dots''), a note in
Danish, and an empty heading ``Sparsity in DML''.

\paragraph{Update, 14 November.} You have since run Chang's nine covariates with
your own first steps (50 splits, $K=5$): $-3.157$ (SE 3.481) with post-Lasso and
$-12.682$ (SE 2.263) without the refit. At your suggestion I then ran Chang's
original replication code (\texttt{Table 2/DMLML.R} in his repository). It
reproduces his published estimate to every digit, and taking it apart changes
what Section~6 should say. Part~\ref{sec:emp-a} is rewritten accordingly. The
scripts behind the numbers are in \texttt{replication/}.

\paragraph{Correction.} The first version of this file said that Chang's
estimate, $-5.222$, equals the unadjusted difference in differences, $-5.228$,
because ``his first steps do essentially no adjustment''. That was wrong. His
Lasso fits select on average 3.7 of the nine covariates for the propensity score
and 2.0 for the outcome; the two numbers agree by coincidence. The claim is
removed from the inserts below and from the conclusion.

The text below is complete except for the numbers in \ph{orange}, which must come
from your run. The numbers in panels A and C of Table~\ref{tab:e0} are from my
runs today; rerun them yourself before you use them (Part~\ref{sec:emp-a} says
how).

Small fixes in the paragraph you keep (``Is trade liberalization a briber's
utopia?''): ``a 8-year window'' should be ``an 8-year window''; and the paragraph
describes the tariff and trade-flow data but not the data Chang uses, the bribe
payments recorded for individual shipments. Add one sentence on where they come
from and how treatment is defined (a tariff change in 2008 for the product), from
Sequeira's paper.

\subsection{Opening, and what the application can show}\label{sec:emp-a}
\where{Replace the opening paragraph of Section~6 (p.~24).}
\begin{ins}
\textcite{chang2020double} illustrates DML-DiD with the study of
\textcite{sequeira2016corruption} on tariff liberalisation and corruption. In
this section I use his application to ask how far the choices of representation
studied in Section~4 move the DML-DiD estimate in real data. The data are
repeated cross-sections, so the estimator is based on the score $\psi_2$.
\end{ins}

\where{After the summary of Sequeira's paper. Replace ``Chang uses only nine
covariates \dots'' and the Danish note.}
\begin{ins}
\paragraph{Chang's specification.} The outcome is $\log(1+\text{bribe})$ for a
shipment, treatment is a change in the tariff of the product in 2008, and the
period indicator is an indicator for shipments after 2008. Chang's replication
code builds dummies for clearing agents and product groups, but the covariate
matrix uses only its first nine columns: the post-2008 indicator, which is the
period indicator $T$ itself; the 2007 tariff; the log value per ton; indicators
for differentiated, agricultural and perishable goods and for large firms; the
day of the week of arrival, entered as a number from 1 to 7; and an indicator
for monitoring. With $N=1084$ complete cases, this is hardly a high-dimensional
setting.

\paragraph{Recreating Chang's estimate.} Chang reports a DML-DiD estimate of
$-5.222$ with a standard error of 2.647. His replication code reproduces both
numbers exactly.\footnote{With R's sampling algorithm set to the one used before
version 3.6.0, \texttt{RNGkind(sample.kind = "Rounding")}. With the current
default, the same code draws different random splits and gives $-5.893$ (SE
2.575).} The code follows Definition~2 with $K=2$ folds, penalties chosen by
10-fold cross-validation with the one-standard-error rule, and no refit, and it
averages the estimate over 100 random splits. In three respects it does
something that the paper does not describe. First, observations with a fitted
propensity score outside $(0.03,0.97)$ are dropped. Second, both nuisance
functions are predicted without the intercept that the Lasso has fitted,
$\hat g=\Lambda(X'\hat\beta)$ and $\hat\ell_2=X'\hat\beta$. The fitted propensity
scores then average 0.27, while 16\% of the shipments are treated. Third, within
each fold, the code discards every score that is at least as large as the
absolute value of the smallest score in the fold before it averages them. This
removes on average 5 of about 530 scores per fold, all of them positive.
\end{ins}
\begin{table}[ht]
\centering\small
\begin{tabular}{l rr rr}
\toprule
 & & & \multicolumn{2}{c}{Single splits}\\
\cmidrule(lr){4-5}
 & Estimate & SE & Min & Max\\
\midrule
\multicolumn{5}{l}{\emph{A. Chang's procedure: cross-validated penalty, no refit, $K=2$, 100 splits}}\\
\quad \textcite[Table~2]{chang2020double} & $-5.222$ & 2.647 & & \\
\quad His replication code & $-5.222$ & 2.647 & $-13.71$ & 7.62\\
\quad \dots\ all scores kept & 4.308 & 4.592 & $-5.35$ & 18.34\\
\quad \dots\ fitted intercepts included & $-12.105$ & 1.711 & $-15.53$ & $-9.17$\\
\quad \dots\ intercepts included and all scores kept & $-10.091$ & 2.377 & $-14.64$ & 0.47\\
\addlinespace
\multicolumn{5}{l}{\emph{B. This thesis: plug-in penalty, clipping, $K=5$, 50 splits}}\\
\quad Lasso (no refit) & $-12.682$ & 2.263 & $-13.69$ & $-11.28$\\
\quad Post-Lasso & $-3.157$ & 3.481 & $-6.15$ & 3.14\\
\addlinespace
\multicolumn{5}{l}{\emph{C. Without first steps}}\\
\quad Difference in differences of the four group means & $-5.228$ & & & \\
\quad Score $\psi_2$, first steps depend on the period only & $-13.33$ & & & \\
\quad Score $\psi_2$, constant first steps & $-18.67$ & & & \\
\bottomrule
\end{tabular}
\caption{DML-DiD in Chang's application with his nine covariates, $N=1084$.
Panel~A: mean over the 100 splits and standard error as computed by Chang's code,
with one step of the code changed at a time. Panel~B: median over the splits and
the standard error of \textcite[Section~3.4]{chernozhukov2018double}.}
\label{tab:e0}
\end{table}
\begin{ins}
Panel~A of Table~\ref{tab:e0} shows what these steps do. The third one decides
the sign: with all scores kept, the same code gives $+4.3$. With the fitted
intercepts included, it gives $-12.1$, and $-10.1$ when all scores are kept as
well. The published estimate thus depends on two steps of the code that the
paper does not describe. Implemented as the paper describes it, with Chang's own
tuning, DML-DiD gives $-10.1$ with a standard error of 2.4, about twice the
published estimate.

Panel~B shows my implementation with the same nine covariates. Without the
refit it gives $-12.7$, close to panel~A with the intercepts included, and
close to what the score gives when the first steps depend on the period only
($-13.3$, panel~C): the shrunken Lasso fits adjust for little more than the
period. With the refit, the estimate is $-3.2$ with a standard error of 3.5, and
single splits range from $-6.1$ to $3.1$.

\paragraph{What the application can show.} For one data set and one set of nine
covariates, the estimates in Table~\ref{tab:e0} range from $+4$ to $-13$. Two
features of the data add to this. First, the estimator for repeated
cross-sections assumes that both periods are samples from the same population
(RCS(i)). In the data, the treated share is 0.40 among the 140 shipments before
2008 and 0.13 among the 944 after, and five of the eight clearing agents appear
only after 2008. This is why the estimate depends so strongly on how the period
indicator, which Chang includes among the covariates, is fitted: with constant
first steps the score gives $-18.7$, and with first steps that depend on the
period only it gives $-13.3$, because the score then weights the post-period
difference between treated and controls by $\Pr(T=1\mid D=1)/\Pr(T=1)=0.78$ and
the pre-period difference by $\Pr(T=0\mid D=1)/\Pr(T=0)=2.46$. Second, 82.3\% of
the outcomes are zero, and an average effect on $\log(1+Y)$ depends on the units
in which $Y$ is measured \parencite{chen2024logs}. I therefore do not interpret
the level of any of these estimates. I keep Chang's specification and use the
application to measure how far the coding of the covariates moves the estimate
within it.
\end{ins}
\note{(i) Panels A and C are my runs. To make the numbers yours: for the second
row, run Chang's \texttt{DMLML.R} unchanged with the line
\texttt{RNGkind(sample.kind = "Rounding")} added before \texttt{DMLML(Y,D,p,T)};
for the other rows of panel~A, use \texttt{replication/chang\_variants.R}
(exact, in R) or the Python port \texttt{replication/chang\_port.py}, which
follows the same steps with other random splits and gives numbers within about
one point of these. (ii) The statement about the paper is a statement about
code, and it is easy to check: the lines are \texttt{beta1hat=fit\$beta},
\texttt{ellhat2=p[index1,]\%*\%beta2hat} and
\texttt{s=s[which(s<abs(min(s)))]} in \texttt{DMLML.R}; neither the paper nor
its online appendix mentions trimming or the removal of scores. Still, show the
table to your supervisor before you build the section on it, and keep the
wording descriptive, as above. (iii) ``for large firms'': the variable is
\texttt{dfs}; check its definition in Sequeira's paper before you describe it.
(iv) $-5.228$ is close to Chang's $-5.222$ by coincidence; do not comment on
it.}

\subsection{Codings and results}\label{sec:emp-b}
\where{After the previous insert.}
\begin{ins}
\paragraph{Codings.} I add the product group of the shipment to Chang's nine
covariates, as dummies. The product groups aggregate the HS 4-digit codes of the
shipments. At the 4-digit level, 107 of the 131 codes in the sample contain only
treated or only untreated shipments, so 4-digit dummies would violate overlap,
and I leave the clearing agents out for the reason given above. One product
group contains 3 shipments, all treated. Without controls, its untreated trend
is not identified, and by the caveat to Lemma~\ref{lem:invariance} even the
unpenalised estimate would depend on the coding, so I drop it ($N=1081$,
13 groups). The 2007 tariff almost determines treatment: products with tariffs
of 0 or 2.5 are never treated, and products with tariffs of 3, 4 or 10 always
are. Dummies or polynomials in the tariff would push the estimated propensity
score to 0 or 1, and I leave it linear. Agricultural products are almost never
treated (2 of 522 shipments), so the comparison effectively rests on
non-agricultural shipments.

I compare three sets of specifications: Chang's nine covariates; the same with
product-group dummies, once for each of the 13 choices of reference group; and
the same with the dummies forced in, which by Lemma~\ref{lem:invariance} does
not depend on the reference group. All estimates use the same \ph{100} random
sample splits with $K=5$, and I report the median over splits with the standard
error of \textcite[Section~3.4]{chernozhukov2018double}, which accounts for the
variation across splits.

\paragraph{Results.} Table~\ref{tab:e1} reports the estimates. The starting point is
the post-Lasso estimate with Chang's nine covariates, $-3.2$ (SE 3.5,
Table~\ref{tab:e0}). With product-group dummies, the
propensity Lasso selects the dummies of \ph{three} groups, \ph{vegetables,
machinery and transportation}, whose treated shares (0.67, 0.67 and 0.05) differ
most from the others. \ph{Nine} of the other ten groups, as reference, yield the
same pools and therefore the same estimate (\ph{$+5.2$}). Only the three groups
whose dummies are selected change the estimate substantially when they serve as
the reference. The default choice in applied work, the largest group, is
transportation, which is atypical in both selection and trends
(Table~\ref{tab:pooling}); it gives the estimate furthest from the benchmark
with the dummies forced in (\ph{$+8.1$} against \ph{$-2.0$}), as
Corollary~\ref{cor:reference} leads one to expect. Across reference groups, the
estimate moves by \ph{12} points, or \ph{1.3} standard errors of the specification
with transportation as the reference.

The standard errors are large, because the 53 treated shipments from before
2008 dominate the score, and the level of the estimate is not credible given the
change in composition between the periods. What the exercise shows is that, in
Chang's application, the choice of reference group alone moves DML-DiD by more
than a standard error, that the sign of the estimate depends on it, and that the
references that matter are the ones Proposition~\ref{prop:pool} points to.
\end{ins}
\note{The orange numbers are from my own run with 20 splits and an earlier
version of the first steps. Replace all of them, and check each sentence against
your table: in particular which dummies are selected, and whether the forced-in
estimate is still negative. If it is not, delete ``that the sign of the estimate
depends on it''. Table~\ref{tab:e1}: one row per specification (Chang's nine
covariates; dummies forced in; each reference group), columns: post-Lasso
estimate (SE), Lasso estimate (SE), number of group dummies selected in
$\hat g$ and in $\hat\ell_2$; below it the range over reference groups in
standard errors.}

\subsection{Centring}
The introduction (p.~2) and the list in Section~4 name three choices of
representation: the reference category, the coding of categorical variables,
and where a variable is centred before it is squared or interacted. The thesis
has theory and simulations for the first two and nothing for the third. Either
add the centring check to Section~6 or remove the third choice from the
introduction and mark it in Section~4 as a choice you do not study. If you run
it, the text is:
\begin{ins}
\paragraph{Centring.} The third choice of Section~4 is where a variable is
centred before it is squared or interacted. I add squares and pairwise
interactions of the continuous covariates to Chang's specification, which gives
a dictionary of \ph{39} columns, and centre the covariates at zero, at their
means or at their medians before forming them. The three dictionaries span the
same space. The estimate moves by \ph{0.24} standard errors across the three
centrings, against \ph{0.06} when only a single square is added.
\textcite{kolesar2026fragility} find ranges above one standard error from
centring, in dictionaries with several hundred columns. In a dictionary of the
size that Chang's sample allows, centring is a second-order concern compared
with the coding of the product groups.
\end{ins}
\note{Linear terms are exactly invariant to centring in your implementation
(standardisation and an unpenalised intercept), so centring can only matter
through the squares and interactions. Say this in one sentence if you include
the paragraph.}

% =============================================================================
\section{Insert: conclusion}\label{sec:concl-ins}

There is no conclusion in v13, although the introduction announces one. The
draft below answers the two questions of the introduction in order and uses
only results that are in the thesis. The paragraph on the application depends on
Section~6.

\where{New Section~7.}
\begin{ins}
\section*{7\quad Conclusion}

This thesis asked how DML-DiD with Lasso first steps behaves when the sparsity
it relies on fails, and which choices of representation make it fail.

The answer to the first question follows from one identity. The bias of DML-DiD
is the product of the errors in the propensity score and in the untreated
outcome trend (Lemma~\ref{lem:bias}). Three consequences follow. First, the
conditions for DML-DiD can be stated for the product: one nuisance may be far
from sparse if the other is sparse enough, which is weaker than Chang's
requirement on each first step, and when the product condition fails, coverage
can go to zero (Propositions~\ref{prop:ratefail} and~\ref{prop:ratedr}). Second,
when the Lasso selects nothing, DML-DiD is the unadjusted DiD, and in general
the bias comes from the covariates that drive both selection and trends and that
both first steps omit (Proposition~\ref{prop:stein}). Third, because the outcome
nuisance is a trend, it does not matter whether the levels of the outcome are
sparse.

The answer to the second question is that the choices that matter are those
that make the Lasso pool units that differ in \emph{both} selection and trends.
With categorical covariates, a post-Lasso fit is constant on groups of
categories, and the bias is the covariance of selection and trends within these
groups (Proposition~\ref{prop:pool}). A reference category that is atypical in
only one of the two is harmless; one that is atypical in both creates a bias in
proportion to its share of the sample (Corollary~\ref{cor:reference}). With
many small categories, a single atypical reference category therefore costs
little. What matters is the size of the category, and the common default of
using the largest category as the reference is the risky one.

The simulations confirm these results and add three observations. The product
condition describes the limit, but close to its boundary the limit is
irrelevant: a design that satisfies it had a coverage of two thirds at every
sample size up to 4{,}000. A bad reference category can leave more bias than no
adjustment at all, and the bias does not fade as the sample grows, so the
estimates under different codings move apart by up to 5.6 standard errors. And
the refit after selection matters throughout: Chang's specification, which plugs
in the penalised fits, had a coverage near zero in designs where post-Lasso was
close to the nominal level. The simulations also show that the design Chang uses
to motivate DML-DiD contains no confounding, so that the unadjusted DiD is
unbiased there.

In Chang's application, \ph{the choice of reference group for the product-group
dummies moves the estimate by 1.3 standard errors, and the default choice, the
largest group, gives the estimate furthest from the invariant benchmark.} The
level of the estimate cannot be interpreted. The composition of the sample
changes between the periods, and for Chang's own nine covariates the estimate
ranges from $+4$ to $-13$ depending on details of the implementation: his
published estimate of $-5.2$ rests on two steps of his code that the paper does
not describe, and without them the same procedure gives $-10.1$.

For applied work, the results suggest four practices. Report DML-DiD under
alternative codings of the controls, together with their range; a large range
indicates that sparsity fails in at least one of them, although a small range
does not prove that it holds. Do not penalise the dummies of a categorical
control when overlap allows it: the estimate is then invariant to the reference
category, at a small cost in precision unless the categories are very small.
Refit after selection. And compare the estimate with the unadjusted DiD: if the
two coincide, the first steps have adjusted for nothing.

The analysis has limits. It covers the two-period, two-group design, and the
theory is worked out for repeated outcomes, with the extension to repeated
cross-sections only sketched. It concerns Lasso first steps; other learners,
such as random forests, are not invariant to the representation either, but the
pooling argument does not apply to them. The simulations use independent
Gaussian covariates and a single categorical covariate. Finally, the thesis
offers no formal test of sparsity for DML-DiD: the tests of
\textcite{kolesar2026fragility} compare with OLS, and an analogue based on the
range across codings, or on the comparison with unpenalised first steps, is a
natural next step, as is the extension to staggered treatment timing.
\end{ins}
\note{``up to 5.6 standard errors'' is the median range in D4a at $N=4000$ and
$\pi_0=0.30$. ``coverage near zero in designs where post-Lasso was close to the
nominal level'' refers to D1 at $N=2000$ (0.000 against 0.946). Adjust the last
paragraph to what you want to claim; ``random forests'' is my addition and not a
result of the thesis, so cut it if you prefer.}

% =============================================================================
\section{Literature}\label{sec:lit}

\subsection{Claims checked against the sources}
I checked the following against the texts of the papers (Chang 2020 and its
appendix; Chernozhukov et al.\ 2018; Kolesár et al., arXiv v5 of 13 May 2026;
Smucler et al.\ 2019; Wüthrich and Zhu 2023; Belloni et al.\ 2012; Belloni and
Chernozhukov 2011; Belloni, Chernozhukov and Wei 2016).

\begin{center}\small
\begin{tabular}{p{0.06\textwidth}p{0.49\textwidth}p{0.36\textwidth}}
\toprule
Page & Claim in v13 & Verdict\\
\midrule
2, 13 & Chang: valid if first steps converge faster than $N^{-1/4}$; Assumption
  3.1 bounds each error by one $\varepsilon_N$ & Correct (his Theorem~3.1 and
  Assumption~3.1).\\
7--9 & Scores, Definition~1, Assumption~3, variance estimators & Agree with
  Chang, apart from the typos in M1--M6. His Definition~3.1 has the
  $\hat p_k$ typo described in M4.\\
10 & Chang chooses the penalty by cross-validation and does not refit & Correct
  (``10-fold cross validations'', p.~189).\\
10 & Plug-in penalty ``Belloni and Chernozhukov (2011)'' & Chang cites
  Belloni et al.\ (2012) for the rule you use. Cite that.\\
10 & Post-Lasso: Belloni et al.\ (2014, JEP) & Wrong source. Belloni and
  Chernozhukov (2013, \emph{Bernoulli}).\\
11 & Kolesár et al.: definition (7) and condition (8) & Correct (their (2) and
  (3)).\\
11 & Kolesár et al., footnote 13: product condition under cross-fitting &
  Correct. The footnote adds that such assumptions are ``likewise sensitive to
  normalization choices'', which is your research question; quote it (M11).\\
4 & Kolesár et al.\ develop two tests & Correct (a residual test and a Hausman
  test, their Section~5).\\
12 & Lemma~4 is the mixed bias property of Smucler et al., eq.~(3) & Correct.
  Their term for the class is ``bilinear influence function''.\\
12 & Assumption~5: Belloni and Chernozhukov (2011, Cor.~1); van de Geer (2008,
  Thm 2.2, Ex.~1); Smucler et al.\ (Thm 9, Remark 7) & Not the estimators you
  use. Replace by M10.\\
13 & ``stricter than necessary (Smucler et al., 2019; Zimmert, 2020)'' & Smucler
  et al.: correct. I could not check Zimmert. The closest source is
  Chernozhukov et al.\ (2018, Remark~5.2); cite it.\\
13--14 & Chernozhukov et al.\ (2018), eq.~(5.4), (3.7), (3.8), Theorem~3.2 &
  Correct. Their Assumption~5.1 for the ATT already has the product-only rate
  condition, and Remark~5.2 states $s^gs^m\ll N$. See Part~\ref{sec:obvious}.\\
5, 15 & Wüthrich and Zhu: under-selection when coefficients are below half the
  penalty; double under-selection & Correct.\\
1 & Snow, 1849 & 1855 (second edition). See Part~\ref{sec:front}.\\
22 & D2 design: Smucler et al., Appendix~D & Correct for $(0.8,1.5)$ and for the
  idea of different decay rates; $(0.7,0.7)$ and $(2,2)$ are your own cells.\\
\bottomrule
\end{tabular}
\end{center}

\paragraph{Not checked.} The session could not reach arXiv, so I could not
check: Rodrigues and Hatfield (2026) and the description ``graphical framework
\dots\ equivalence between parallel trends and equi-confounding''; the NBER
number of Goldsmith-Pinkham (2026); Hünermund et al.\ (2023); Zimmert (2020),
described as ``builds on the same setup'' (the paper is independent of Chang's
and contemporaneous, so ``considers the same problem'' is safer); and whether
Callaway et al.\ (2025), Caetano and Callaway (2026) and Karim and Webb (2026)
have been published in a journal in the meantime. Check these yourself.

\subsection{References to add}
\begin{itemize}
\item \textcite{belloni2012sparse}: the penalty loadings (Section~5) and the rate
  of the outcome Lasso (Assumption~5).
\item \textcite{belloni2013least}: post-Lasso.
\item \textcite{belloni2016post}: logit Lasso, its penalty and its rate.
\item \textcite{belloni2014inference}: post-double-selection, the estimator
  Kolesár et al.\ study; also polynomially decaying designs.
\item \textcite{farrell2015robust}: the product condition for doubly robust
  estimators.
\item \textcite{belloni2017program} and \textcite{athey2018approximate}:
  polynomially decaying designs (D2/D3).
\item \textcite{hirano2003efficient}: estimated versus true propensity score
  (D1, footnote).
\item \textcite{sur2019modern}: bias of the unpenalised logit when $p/n$ is not
  small (caveat to Lemma~10).
\item \textcite{chen2024logs}: $\log(1+Y)$ (Section~6).
\item \textcite{snow1855mode} and \textcite{angrist2009mostly}.
\end{itemize}
The entries are in \texttt{review/v13\_refs.bib}.

\subsection{The reference list as printed (pp.~25--27)}
\begin{itemize}
\item ``Angrist, J. D. (n.d.)'': Angrist and Pischke (2009), Princeton University
  Press.
\item ``Van der Vaart, A. (n.d.). \emph{Asymptotic statistics} (Vol. 1998)'': year
  1998, and the DOI is printed twice in the link.
\item ``Belloni, A., \& Chernozhukov, V. (2011, September 2)'': it is a chapter in
  \emph{Inverse Problems and High-Dimensional Estimation} (Springer, 2011);
  cite the chapter, not the arXiv date.
\item Smucler et al.: the title prints ``\verb|$\ell_1$|'' as raw LaTeX.
\item Sequeira: ``southern africa'' needs capitals (protect them with braces in
  the \texttt{.bib}).
\item ``Van De Geer'' should be ``van de Geer''.
\item ``(Eilers et al.)'' in Definitions~21--22 has no year and no entry; Pitman
  (1993) is cited only for Markov and Chebyshev. Both disappear if you cut that
  material (Part~\ref{sec:obvious}).
\item Belloni et al.\ (2014, JEP) is no longer cited once the post-Lasso
  reference is corrected.
\end{itemize}

% =============================================================================
\section{Stating the obvious? What is new and what is not}\label{sec:obvious}

You asked whether the thesis only states the obvious. It does not, but it
distributes its emphasis the wrong way round: the known results are presented
as contributions and proved at length, and the original ones are not mentioned
in the introduction.

\begin{center}\small
\begin{tabular}{p{0.20\textwidth}p{0.36\textwidth}p{0.35\textwidth}}
\toprule
Result & Status & What to do\\
\midrule
Lemma 4 (mixed bias) & Known (Smucler et al.; implicit in Chernozhukov et al.). &
  Keep, as now, with the attribution. Add the proof (M9).\\
Proposition 6 (product of sparsity indices) & Follows in one line from
  Assumption~5. Known: Chernozhukov et al.\ (2018, Remark~5.2), Farrell (2015),
  Kolesár et al., footnote~13. & Keep, but say that it is known (M11). Delete
  the two sentences after it (``Both parts follow directly \dots\ goes to zero
  exactly when $s\log p/N$ does''), which restate it.\\
Proposition 8 (rate double robustness) & Theorem~5.1 of Chernozhukov et al.\
  applied to the change in outcomes. & Present it as that (M13). Its point is
  that Chang's condition is stronger than needed.\\
Proposition 7 (rate failures) & (a) is routine. (b)--(c), an explicit sequence
  for which coverage goes to zero, I have not seen written down; Chernozhukov et
  al.\ assert sharpness without proof. & Yours. Keep. Say that (b) is a worst
  case and link it to D2/D3 (M12).\\
Remark (no adjustment) & Elementary. & Keep as a remark. Its value is the
  benchmark numbers it gives for the simulations (1.563, $0.826\,\gamma'\delta$,
  0.184).\\
Proposition 9 (Gaussian bias) & Yours. The ``both omit'' form is the DML-DiD
  counterpart of Wüthrich and Zhu. & Keep and emphasise. Use the simpler
  constants (M15).\\
Lemma 10 (invariance) & Elementary; Kolesár et al.\ note it for OLS. & Keep,
  short. The caveats (separation, categories without controls) are the part
  that is not obvious and that Section~6 needs.\\
Lemma 11, Proposition 12 & Algebra. The decomposition is never used again: no
  table reports its three terms and the bound is never evaluated. & See below.\\
Proposition 13, Corollary 14 (pooling, reference category) & Yours, and the most
  original part of the thesis. & Put them in the introduction (Part~\ref{sec:front}).\\
D1 & Partly a replication. New: Chang's design has no confounding; why Abadie's
  estimator is biased there; the role of the refit. & Say these three things
  explicitly (Part~\ref{sec:sims}).\\
D2--D4 & Yours. & Part~\ref{sec:sims}.\\
\bottomrule
\end{tabular}
\end{center}

\paragraph{Proposition 12.} Two options. (i) Reduce it: state the identity in
one display in the text, keep the observation that treated observations are
unaffected by the propensity fit (so that a change of coding acts through the
weights of the controls and through the outcome fit), keep the definition of
$\mathcal R(\mathcal A)$, move Lemma~11 to the appendix, where the proof of
Lemma~27 uses it, and drop the bound. (ii) Use it: report the three terms for
one pair of codings in D4 or in Section~6. I recommend (i); it saves a page in
the text and a page and a half in the appendix.

\paragraph{The appendix.} The first four pages (pp.~28--32) are textbook
material, and most of it is not used in any proof.
\begin{itemize}
\item Cut Definition~15 (limits), the empty Definition~16, and Definitions~20--22
  (epsilon functions, $o$ and $O$ for functions). Nothing refers to them.
\item Cut the proof of Chebyshev's inequality from Markov's (Definitions~17--18,
  Lemma~19). You use Chebyshev once, for $O_k=O_P(n^{-1/2})$. Write ``by
  Chebyshev's inequality'' there.
\item Keep the $o_P$/$O_P$ notation (Definition~23) as a short paragraph, and
  Definition~24 as a lemma with its source (Chernozhukov et al.\ 2018,
  Lemma~6.1). Markov, Chebyshev and Definition~24 are results, not definitions.
\item ``Proof of zero mean adjustment term'' is two lines. Put it in
  Section~3.2 as one sentence: $\E[c_j\mid X]=0$ because $\E[D\mid X]=g_0$.
\item Cut the ``Walk-through of lasso estimation'': it derives the logistic
  log-likelihood, which is already displayed in Section~3.3.
\item Cut the paragraph that apologises to the reader (p.~28).
\end{itemize}
This removes about four pages and loses nothing an examiner would look for.

\paragraph{In the text.} Sentences that tell the reader nothing can go: the
second half of the paragraph after Proposition~6; ``This is the upper bound of
the error in estimation, and the realized error can be much smaller'' (p.~38);
``$g_N$ is thus a legitimate probability'' (p.~38) can be one clause.

% =============================================================================
\section{Coherence}\label{sec:coherence}

\begin{enumerate}
\item \textbf{The two questions of the introduction are never answered in one
  place.} The conclusion in Part~\ref{sec:concl-ins} does that.
\item \textbf{The two kinds of failure.} The introduction distinguishes failure
  (1), sparse in some basis but not in the researcher's, and failure (2), not
  sparse in any basis. Section~4.1 is about (2) and Section~4.2 about (1); D2/D3
  test (2) and D4 tests (1). Say so in the first sentence of each of these
  parts. The inserts for Section~5 do it for the designs.
\item \textbf{Section~4 opens with two other ``questions''} (p.~10: ``The first
  question to emerge is whether sparsity applies to DML-DiD \dots\ The second
  question to emergu \dots''). They are not the research questions, and the
  first sentence is ungrammatical. Suggested replacement:
\begin{ins}
In this section, I bring the critique of \textcite{kolesar2026fragility} to the
estimator of the previous section. DML-DiD itself does not assume sparsity: it
needs first-step estimators that converge fast enough. Sparsity enters through
the choice of a sparsity-based first step, such as the Lasso, and DML-DiD is
attractive precisely because it appears to leave the choice of covariates to
the first step. But a Lasso depends on how the covariates are written down, and
\textcite{kolesar2026fragility} document the following choices, none of which
affects OLS. I state all results for repeated outcomes; they carry over to
repeated cross-sections with $\psi_2$ and $\ell_{20}$, as noted after
Lemma~\ref{lem:bias} and at the end of the proof of Proposition~\ref{prop:ratefail}.
\end{ins}
\item \textbf{Repeated outcomes and repeated cross-sections.} The theory and all
  simulations are for repeated outcomes; the application uses repeated
  cross-sections. The last sentence of the insert above and the first paragraph
  of Section~6 make the link.
\item \textbf{Lasso and post-Lasso.} Section~3.3 says you refit, Assumption~5
  speaks of ``lasso first steps'', Section~5 reports both. Write ``Lasso or
  post-Lasso first steps'' in Assumption~5, and call the two versions the same
  thing everywhere: ``post-Lasso'' and ``Lasso (no refit)''.
\item \textbf{$p\asymp N$.} Section~4 assumes $p\asymp N$. D1--D3 have $p=N/2$.
  D4 has a fixed number of categories, so asymptotically the Lasso would select
  every relevant dummy; D4 is about what happens before that. Result~3 of D4
  says so; keep that sentence.
\item \textbf{Table~1 does not match the list of estimators}, and D3 is
  undefined. Fixed by the inserts.
\item \textbf{Headings.} ``Priors'' has a Bayesian meaning; ``Predictions'' is the
  usual word. ``Digest'' could be ``What D1 shows''. I kept ``Priors'' in the
  inserts to match your headings; change it if you agree.
\item \textbf{The Remark has no number}, which is why ``as per~.'' is empty.
\item \textbf{There is no abstract} in v13. Check what your study programme
  requires (abstract, front page, word count).
\end{enumerate}

\paragraph{Notation used for more than one thing.}
\begin{center}\small
\begin{tabular}{p{0.10\textwidth}p{0.52\textwidth}p{0.30\textwidth}}
\toprule
Symbol & Meanings in v13 & Suggestion\\
\midrule
$c$, $c_j$ & adjustment terms $c_1,c_2$ in the scores (p.~7); constants
  $c_1,c_2$ in Proposition~9; terms $c_1,c_2,c_3$ in Proposition~12; rate
  exponent $c$ in Proposition~7; penalty constant $c=1.1$; constant fit
  $\hat\ell\equiv c$; category $c$; coverage $c$ & $\bar c_1,\bar c_2$ in
  Proposition~9; (I)--(III) in Proposition~12; exponents $b_g,b_\ell$; coverage
  in words\\
$\delta$ & trend coefficients; $\delta_g=g'-g$ (Lemma~11); $\delta_g,\delta_\ell$
  errors (Lemma~27); a constant in the proof of Proposition~7(c) & $\Delta g$ in
  Lemma~11; another letter in the proof\\
$\gamma$ & propensity coefficients; penalty loadings $\gamma_{kj}$ (p.~10);
  Chang's level $\gamma\to0$ (p.~19) & $\hat\psi_{kj}$ for the loadings\\
$\alpha$ & level of the penalty; level of the test; decay rates
  $\alpha_g,\alpha_\ell$ & acceptable; say ``decay rates'' each time\\
$n$ & fold size; sample size in the Lasso (p.~19) & $M$ in the Lasso\\
$s$ & dimension of $\theta_0$ (p.~7); sparsity index & drop the first\\
$w$ & argument of $\Lambda$; $w(X)$ and its mean (Prop.~9); pool weight
  (Cor.~14); $w_g,w_0$ (Lemma~27); $w_i^{(A)}$ (p.~42) & $\bar w$ for the mean;
  $\omega$ in Corollary~14\\
$\tilde\theta$, $\hat\theta$ & both used for the final estimator & $\tilde\theta$
  (M5)\\
$\lambda$ & $P(T=1)$; penalty level & acceptable, already flagged on p.~10\\
\bottomrule
\end{tabular}
\end{center}

% =============================================================================
\section{Form}\label{sec:form}

\paragraph{Unresolved references.} ``Lemma ??'' on p.~20 and in the caption of
Table~1 (Lemma~4); p.~36 and p.~42 (Lemma~11); ``corollary ??'' and ``Lemma ??''
on pp.~23--24 (Corollary~14, Lemma~10); ``as per~.'' on p.~20 (the Remark); ``6
below'' on p.~12; ``(9)'' and the red ``repeated outcomes'' on p.~7; ``27(ii)''
on p.~13. In the appendix, ``Lemma (4)'' and ``inequality (17)'', ``(19)'' print
like equation numbers.

\paragraph{Broken citations.} Listed in Part~\ref{sec:front}.

\paragraph{Left-over notes.} ``INDSÆT HER'' (p.~17); the stray ``I'' (p.~24); the
Danish note and the heading ``Sparsity in DML'' (p.~25); the empty
Definition~16.

\paragraph{Typos I noticed} (run a spell-check as well):
p.~1 ``treated an untreated'';
p.~2 ``ML method innocuous tools'', ``choices that does not'';
p.~3 ``one nuisances'', ``chang uses'';
p.~4 ``commmon'', ``simultanously'', ``number of covariates exceed'';
p.~5 ``two-nusaince'', ``both panel and repeated outcomes'' (repeated outcomes
and repeated cross-sections);
p.~7 ``$j=\{1,2\}$'' ($j\in\{1,2\}$);
p.~8 ``nuissance'';
p.~9 ``of of'';
p.~10 ``of he nuisance'', ``I my thesis'', ``the size of the set where'';
p.~11 ``Since DML-DiD that allows'', ``a sparsity-based estimators'',
``emergu'';
p.~13 ``suffiecient'';
p.~14 ``uandjusted'', ``sparse in large'';
p.~15 ``that escape both fits omit'';
p.~16 ``Anecdotically'', ``$\hat g_K$'';
p.~18 ``were selection'' (where), ``If it atypical'', ``intutive'', ``leave the
coefficients'', ``neglectible'', ``standard erro'', ``treshhold'';
p.~19 ``lasso ust'', ``samling'', ``Undjusted'';
p.~20 ``be follow'', ``maintaining letting'';
p.~22 ``requieres'', ``how DML-DiD response'';
p.~23 ``attatched'', ``suffiecient'', ``the how'';
p.~24 ``selct'';
p.~25 ``a 8-year'', ``commmon'', ``used create'', ``conflate the covariate
matrix'' (you mean ``expand''), and ``Since squares and interactions of dummy
variables are degenerate, I will include them'' (you mean ``exclude'');
appendix: ``supress'', ``Walk-trough'', ``estimation of of'', ``asusmption'',
``$q_pX_i$'', ``th whole'', ``$\hat\lambda_K$''.
Also write ``Lasso'' with a capital throughout (it is an acronym), and choose
between British and American spelling (``categorise''/``categorize'',
``centred''/``centered'' both occur).

% =============================================================================
\section{Questions for you}\label{sec:questions}

\begin{enumerate}
\item \textbf{Section~6.} The baseline and the recreation of Chang's estimate
  are done (Table~\ref{tab:e0}). Do you want the full comparison across the 13
  reference groups and the forced-in dummies (about $15\times50$ cross-fits), or
  two or three reference groups only? And are you comfortable reporting panel~A
  of Table~\ref{tab:e0}? Ask your supervisor.
\item \textbf{Centring.} Include the check in Section~6, or drop centring as a
  choice you study?
\item \textbf{D2/D3.} Main text with both tables, or one table in the text and
  the rest in the appendix? Is there a page limit?
\item \textbf{Titles.} Has your supervisor agreed to the informal title and
  section titles?
\item \textbf{Forced-in dummies.} Do you want to show the invariance to the
  reference category numerically (one extra run), or rely on Lemma~10?
\item \textbf{Deadline.} If you tell me the date, I can order the remaining work
  by what an examiner will weigh most.
\end{enumerate}

\printbibliography[title={References used in the inserts}]
\end{document}
